\chapter{Geometria Kleina: przestrzeń i jej symetrie}
\label{ch:geometria-kleina}
\index{geometria Kleina}
\index{przestrzeń jednorodna}

Grupa Liego z Rozdziału~\ref{ch:grupy-liego} może działać na innej
rozmaitości. Zamiast pytać tylko o punkty i współrzędne, możemy
zapytać, które własności punktów, prostych lub odległości nie zmieniają
się pod tym działaniem. Taki był punkt wyjścia programu erlangeńskiego
Kleina: geometrię opisuje przestrzeń wraz z wybraną grupą jej
przekształceń. Ta sama rozmaitość może więc nieść różne geometrie.

W tym rozdziale zbudujemy język przestrzeni jednorodnych $G/H$.
Przykłady prowadzą od płaszczyzny euklidesowej i afinicznej przez
przestrzeń rzutową i sferę do torusa oraz dysku Poincarégo. Metryki
hiperboliczne wprowadziliśmy już w
Przykładzie~\ref{prz:trzy-modele-hiperboliczne}; tutaj zbadamy działanie
grupy $\mathrm{PSL}(2,\R)$ przez izometrie w wymiarze dwa.

\section{Działanie, orbita i stabilizator}
\label{sec:klein-dzialanie}

\begin{definicja}[Gładkie działanie grupy Liego]
\label{def:klein-dzialanie}
\index{działanie grupy Liego}
Niech $G$ będzie grupą Liego, a $M$ gładką rozmaitością bez brzegu.
\emph{Gładkim działaniem lewym} $G$ na $M$ nazywamy gładkie
odwzorowanie
\[
 a:G\times M\longrightarrow M,\qquad a(g,p)=g\cdot p,
\]
dla którego $e\cdot p=p$ oraz $(gh)\cdot p=g\cdot(h\cdot p)$.
Każdy $g$ daje dyfeomorfizm $a_g:p\mapsto g\cdot p$, którego
odwrotnością jest $a_{g^{-1}}$, ponieważ
$a_g\circ a_{g^{-1}}=a_{gg^{-1}}=a_e=\id_M$
i tak samo w drugiej kolejności. Intuicyjnie grupa jest zbiorem
dozwolonych ruchów przestrzeni. Na przykład grupa $(\R^n,+)$ działa
na $\R^n$ przez translacje $b\cdot x=x+b$.
\end{definicja}

\begin{definicja}[Orbita, stabilizator i skuteczność]
\label{def:klein-orbita}
\index{orbita działania}
\index{stabilizator punktu}
\index{działanie skuteczne}
Dla $p\in M$ określamy \emph{orbitę} i \emph{stabilizator}
\[
 G\cdot p=\{g\cdot p:g\in G\},\qquad
 G_p=\{g\in G:g\cdot p=p\}.
\]
Działanie jest \emph{przechodnie}, gdy $G\cdot p=M$ dla pewnego,
a więc każdego $p\in M$. Jest \emph{skuteczne}, gdy jedynym elementem
ustalającym każdy punkt jest $e$. Orbita mówi, dokąd można przesunąć
punkt; stabilizator mówi, które ruchy pozostawiają go na miejscu.
Przy translacjach $\R^n$ na $\R^n$ orbita każdego punktu jest całą
przestrzenią, a stabilizator jest trywialny.
\end{definicja}

Stabilizator jest podgrupą: $e\in G_p$, a jeśli $g,h\in G_p$, to
$gh^{-1}\cdot p=g\cdot(h^{-1}\cdot p)=p$. Jest też domknięty,
gdyż jest przeciwobrazem punktu $p$ przez ciągłe odwzorowanie
$g\mapsto g\cdot p$, a rozmaitość $M$ jest Hausdorffa. Nie każda
orbita musi być domknięta ani być podrozmaitością z topologią
odziedziczoną z $M$; to rozróżnienie będzie potrzebne w twierdzeniu
poniżej.

\begin{lemat}[Stały rząd odwzorowania orbity]
\label{lem:klein-rzad-orbity}
Niech $p\in M$ i $o_p:G\to M$, $o_p(g)=g\cdot p$. Rząd
$\dd(o_p)_g$ nie zależy od $g$. Jego jądro w $e$ jest
$T_eG_p$, a $G_p$ jest zanurzoną podrozmaitością $G$ wymiaru
$\dim G-\operatorname{rank}\dd(o_p)_e$.
\end{lemat}
\begin{proof}
Dla $g,h\in G$ zachodzi $o_p\circ L_g=a_g\circ o_p$.
Różniczkując tę równość w $h$ i pamiętając, że $L_g$ oraz
$a_g$ są dyfeomorfizmami, dostajemy
\[
 \dd(o_p)_{gh}\circ\dd(L_g)_h
 =\dd(a_g)_{h\cdot p}\circ\dd(o_p)_h.
\]
Obie skrajne różniczki, $\dd(L_g)_h$ i $\dd(a_g)_{h\cdot p}$,
są izomorfizmami. Mnożenie odwzorowania liniowego przez
izomorfizmy nie zmienia jego rzędu, więc
$\operatorname{rank}\dd(o_p)_{gh}
=\operatorname{rank}\dd(o_p)_h$.
Odwzorowanie $o_p$ ma więc stały rząd $r$. Z
Twierdzenia~\ref{tw:rzad-staly} w mapach wokół $h\in G_p$ i $p$
ma ono postać $(u,v)\mapsto(u,0)$, gdzie $u\in\R^r$.
Jego przeciwobraz punktu $p$ jest lokalnie zadany równaniem $u=0$.
To dowodzi, że $G_p=o_p^{-1}(p)$ jest zanurzoną podrozmaitością
o podanym wymiarze i że $T_eG_p=\ker\dd(o_p)_e$.
Ograniczenia mnożenia i odwrotności w $G$ są gładkie, zatem
$G_p$ jest także grupą Liego z otrzymaną strukturą.
\end{proof}

Do opisu $G/H$ dla \emph{dowolnej} domkniętej podgrupy $H$ użyjemy
standardowego wyniku o grupach Liego. Jego założenie domkniętości
jest istotne: gęsta właściwa podgrupa może dać iloraz, w którym
nie da się rozdzielić dwóch klas otwartymi zbiorami.

\begin{twierdzenie}[Iloraz przez domkniętą podgrupę; wynik zewnętrzny]
\label{thm:klein-iloraz-liego}
Niech $G$ będzie skończenie wymiarową grupą Liego, a $H\subseteq G$
domkniętą podgrupą. Wtedy $H$ ma jednoznaczną strukturę zanurzonej
podgrupy Liego. Zbiór lewych warstw $G/H=\{gH:g\in G\}$ z topologią
ilorazową ma jednoznaczną strukturę gładkiej rozmaitości, dla której
rzutowanie $q:G\to G/H$ jest submersją. Zachodzi
\[
 \dim(G/H)=\dim G-\dim H,\qquad
 T_{eH}(G/H)\cong T_eG/T_eH.
\]
Działanie $G$ na $G/H$ przez $a\cdot(gH)=(ag)H$ jest gładkie
i przechodnie.
\end{twierdzenie}

Jest to połączenie twierdzenia o domkniętej podgrupie z twierdzeniem
o ilorazie przez wolne działanie grupy Liego spełniające warunek
właściwości z Definicji~\ref{def:klein-dzialanie-wlasciwe}; dokładne
wersje i dowody podaje J.\ M.\ Lee, \emph{Introduction to Smooth
Manifolds}, rozdziały~20--21. W drugim twierdzeniu $H$ działa na
$G$ z prawej strony przez $(g,h)\mapsto gh$. Działanie jest wolne.
Aby sprawdzić jego właściwość, niech $C\subset G\times G$ będzie
zwarty i niech $K_1,K_2$ będą jego rzutami na czynniki.
Przeciwobraz $C$ przez $(g,h)\mapsto(g,gh)$ jest domknięty
w zwartym zbiorze
$K_1\times(H\cap K_1^{-1}K_2)$: drugi czynnik jest zwarty,
bo $K_1^{-1}K_2$ jest zwarty, a $H$ domknięta.
Zatem ten przeciwobraz jest zwarty. Wzór na wymiar i przestrzeń
styczną wynika już z lokalnej postaci submersji $q$ z
Wniosku~\ref{wn:submersja-lokalnie}, ponieważ
$q^{-1}(eH)=H$.

\begin{twierdzenie}[Orbita jako przestrzeń jednorodna]
\label{thm:klein-orbita-iloraz}
Niech grupa Liego $G$ działa gładko na rozmaitości $M$ bez brzegu
i niech $p\in M$, $H=G_p$. Wówczas
\[
 F:G/H\longrightarrow M,\qquad F(gH)=g\cdot p,
\]
jest dobrze określoną, różnowartościową immersją o obrazie
$G\cdot p$. Orbita ma przez $F$ kanoniczną strukturę gładkiej
rozmaitości z immersją w $M$, nawet jeśli jej topologia jako
podzbioru $M$ jest inna. Jeśli działanie jest przechodnie,
$F:G/H\to M$ jest
dyfeomorfizmem. Jeśli dana orbita jest zanurzoną podrozmaitością
$M$, to $F$ jest dyfeomorfizmem również na tę podrozmaitość.
\end{twierdzenie}
\begin{proof}
Z Lematu~\ref{lem:klein-rzad-orbity} stabilizator jest domkniętą
zanurzoną podgrupą, więc $G/H$ istnieje na mocy
Twierdzenia~\ref{thm:klein-iloraz-liego}. Równość
$gH=g'H$ oznacza $g'=gh$ dla pewnego $h\in H$; wtedy
$g'\cdot p=g\cdot(h\cdot p)=g\cdot p$. Odwrotnie, jeśli
$g'\cdot p=g\cdot p$, to $g^{-1}g'\in H$, zatem $g'H=gH$.
To daje poprawność definicji i różnowartościowość $F$.

Rzutowanie $q$ ma lokalne gładkie przekroje dzięki postaci
submersji. Ponieważ $F\circ q=o_p$, złożenie $o_p$ z takim
przekrojem $s:U\to G$ daje
$F|_U=F\circ q\circ s=o_p\circ s$, a więc gładkość $F$.
Różniczkując w $e$, otrzymujemy
\[
 \dd F_{eH}\circ\dd q_e=\dd(o_p)_e.
\]
Jądra $\dd q_e$ i $\dd(o_p)_e$ są równe $T_eH$;
stąd $\dd F_{eH}$ jest różnowartościowe.
Dokładniej, jeśli $\dd F_{eH}(v)=0$, wybieramy $X\in T_eG$
z $\dd q_e(X)=v$, co jest możliwe dzięki surjektywności
$\dd q_e$. Wtedy $\dd(o_p)_e(X)=0$, zatem
$X\in T_eH=\ker\dd q_e$ i $v=0$.
Ta sama własność
w każdym $gH$ wynika z $G$-równoważności $F$ i z tego, że
działania przez $g$ na obu rozmaitościach są dyfeomorfizmami.
Zatem $F$ jest immersją.

Jeśli $G\cdot p=M$ i $\dim M>0$, stały rząd $o_p$ musi
równać się $\dim M$. Gdyby był mniejszy, każdy punkt $G$
byłby krytyczny, a Twierdzenie Sarda~\ref{thm:sard} mówiłoby,
że obraz $o_p(G)=M$ ma miarę zero w $M$. To niemożliwe,
bo niepusty otwarty zbiór w mapie ma dodatnią miarę.
Dla $\dim M=0$ wniosek o submersji jest natychmiastowy.
Zatem $\dd(o_p)_e$ jest surjektywne, a powyższa równość
różniczek pokazuje, że $\dd F_{eH}$, a więc każda
$\dd F_{gH}$, jest izomorfizmem. Twierdzenie o funkcji
odwrotnej czyni $F$ lokalnym dyfeomorfizmem. Jest ono także
bijekcją na $M$, więc jego lokalne odwrotności składają się
na gładką odwrotność globalną.

Jeśli zamiast $M$ rozważymy zanurzoną orbitę $N=G\cdot p$,
odwzorowanie $G\times N\to N$ jest gładkie:
w mapie dostosowanej do podrozmaitości $N$ współrzędne
poprzeczne obrazu są stale zerowe, a pozostałe są
ograniczeniami gładkiego działania na $M$.
Tak samo $F:G/H\to N$ jest gładkie. Jego różniczki mają ten sam rząd
$\dim(G/H)$, zaś argument Sarda zastosowany do przechodniego
działania na $N$ daje $\dim N=\dim(G/H)$. Stosujemy więc
poprzedni przypadek do $N$.
\end{proof}

Ostatnie założenie w twierdzeniu ma znaczenie. Działanie
$\R$ na $T^2$ przez
$t\cdot[x,y]=[x+t,y+\alpha t]$ dla niewymiernego $\alpha$
ma trywialny stabilizator, lecz orbita $[0,0]$ jest gęsta
i nie jest podrozmaitością z topologią odziedziczoną z torusa.
Jest to nawinięcie niewymierne opisane po
Przykładzie~\ref{prz:lie-torus-podgrupy}; odwzorowanie orbity jest
różnowartościową immersją, ale nie zanurzeniem w znaczeniu
Rozdziału~\ref{ch:rozmaitosci}.

\begin{definicja}[Geometria Kleina]
\label{def:klein-geometria}
\index{geometria Kleina!definicja}
\emph{Geometrią Kleina} będziemy nazywać parę $(G,X)$, w której
$X$ jest spójną gładką rozmaitością, a $G$ grupą Liego działającą
na niej gładko i przechodnio. Jej \emph{model skuteczny} ma ponadto
skuteczne działanie. Wybór $x_0\in X$ przedstawia przestrzeń jako
$X\cong G/G_{x_0}$; stabilizator zachowuje wybrany punkt.
Na przykład $(\R^n,\R^n)$ z translacjami jest geometrią Kleina,
w której stabilizator zera jest trywialny.
\end{definicja}

Niektórzy autorzy wymagają skuteczności już w definicji. Możemy
zawsze usunąć nieskuteczność: jądro działania jest
\begin{equation}\label{eq:klein-jadro}
 K=\{g:g\cdot x=x\text{ dla każdego }x\in X\}
   =\bigcap_{g\in G}gHg^{-1},\qquad H=G_{x_0}.
\end{equation}
Druga równość wynika z $x=g\cdot x_0$: element $a$ ustala
$x$ dokładnie wtedy, gdy $g^{-1}ag\in H$. Jest to domknięta
podgrupa normalna. Dlatego $G/K$ jest grupą Liego i działa
skutecznie oraz przechodnio na $X$. W szczególności stabilizator
$H$ nie musi być trywialny, aby całe działanie było skuteczne.

Warto zobaczyć różnicę na płaszczyźnie.
Ruchy $g(x)=x+b$ i $g'(x)=R_\theta x+b$, gdzie $R_\theta$
jest obrotem wokół zera, wysyłają $0$ w ten sam punkt $b$.
Różnią się jednak na pozostałych punktach.
Dla grupy ruchów euklidesowych
$E(2)=\{x\mapsto Ax+b:A\in\mathrm O(2),\ b\in\R^2\}$
stabilizator zera to właśnie $\mathrm O(2)$;
pełny rachunek podamy w Przykładzie~\ref{prz:klein-euklidesowa}.
Zapis $E(2)/\mathrm O(2)$ zapomina o wyborze
takiego obrotu lub odbicia przed przesunięciem:
cała warstwa $g\mathrm O(2)$ reprezentuje jeden punkt $b$.

\begin{figure}[htbp]
\centering
\begin{tikzpicture}[>=Stealth,scale=.9]
 \fill[manifoldblue!6] (-.9,-.75) rectangle (5.1,2.6);
 \draw[->,manifoldblue!40] (-.7,0)--(4.9,0);
 \draw[->,manifoldblue!40] (0,-.6)--(0,2.4);
 \coordinate (O) at (0,0);
 \coordinate (B) at (3.4,1.2);
 \draw[->,thick,manifoldblue] (O)--(B)
  node[midway,below=4pt] {$g$};
 \draw[->,thick,manifoldgold!90!black] (O) to[bend left=28]
  node[midway,above=4pt] {$g'=g\circ R_\theta$} (B);
 \draw[->,manifoldred,thick] (.65,0) arc[start angle=0,end angle=135,radius=.65];
 \node[manifoldred] at (-.45,1.05) {$R_\theta$};
 \draw[->,thick,manifoldblue] (B)--++(.85,0);
 \draw[->,thick,manifoldgold!90!black] (B)--++(-.6,.6);
 \fill[manifoldblue!80!black] (O) circle(2.1pt) (B) circle(2.1pt);
 \node[below left=3pt] at (O) {$0$};
 \node[below right=4pt] at (B) {$b$};
 \node[anchor=west,align=left,font=\small,text width=5cm] at (6,1)
  {Stabilizator $\mathrm O(2)$ ustala $0$.\\[5pt]
   Różne ruchy w warstwie $g\mathrm O(2)$ dają ten sam punkt $b$,
   ale różnie przenoszą kierunki.};
\end{tikzpicture}
\caption{Punkt przestrzeni jednorodnej jako warstwa grupy.
Strzałki przy $b$ pokazują obrazy tego samego kierunku:
translacja i translacja po obrocie przenoszą go inaczej.}
\label{fig:klein-orbita-stabilizator}
\end{figure}

\section{Metryka a wybór grupy}
\label{sec:klein-metryka}

Struktura $G/H$ sama nie wybiera metryki. Można to uściślić,
korzystając z działania adjungowanego z
Definicji~\ref{def:lie-ad}.

\begin{stwierdzenie}[Kiedy istnieje metryka niezmiennicza]
\label{prop:klein-metryka-niezmiennicza}
Niech $H$ będzie domkniętą podgrupą $G$, a
$\mathfrak g=T_eG$, $\mathfrak h=T_eH$. Metryki riemannowskie
na $G/H$ niezmiennicze względem działania $G$ odpowiadają
dodatnio określonym iloczynom skalarnym na
$\mathfrak g/\mathfrak h$, niezmienniczym względem
indukowanego działania $\operatorname{Ad}(H)$.
\end{stwierdzenie}
\begin{proof}
Każda metryka jest wyznaczona przez swoje wartości w punktach.
Jeśli jest $G$-niezmiennicza, jej wartość w $eH$ musi być
niezmiennicza pod stabilizatorem $H$. Utożsamienie
$T_{eH}(G/H)\cong\mathfrak g/\mathfrak h$ z
Twierdzenia~\ref{thm:klein-iloraz-liego} przenosi działanie
$h\in H$ na $\operatorname{Ad}_h$ modulo $\mathfrak h$.
Jest ono dobrze określone, bo sprzężenie przez $h$
zachowuje $H$, więc $\operatorname{Ad}_h\mathfrak h=\mathfrak h$.
Istotnie,
\[
 h\exp_G(tX)H
 =h\exp_G(tX)h^{-1}H
 =\exp_G(t\operatorname{Ad}_hX)H.
\]
Różniczkowanie w $t=0$ daje tezę o reprezentacji na
przestrzeni stycznej.

Odwrotnie, oznaczmy dany iloczyn skalarny w $eH$ przez $b$.
Dla $u,v\in T_{gH}(G/H)$ definiujemy
\[
 g_{gH}(u,v)
 =b\bigl(\dd(a_{g^{-1}})_{gH}u,\dd(a_{g^{-1}})_{gH}v\bigr).
\]
To jest dokładny sens przeniesienia pomiaru z $eH$ do $gH$.
Jeśli $gH=g'H$ z $g'=gh$, dwa sposoby przeniesienia
różnią się działaniem $h^{-1}$ w $eH$:
$a_{(g')^{-1}}=a_{h^{-1}}\circ a_{g^{-1}}$.
Różniczka pierwszego czynnika zachowuje $b$,
więc oba wzory dają tę samą wartość.
Dla $k\in G$ równość $(kg)^{-1}k=g^{-1}$
i reguła łańcuchowa dają
$g_{kgH}(\dd(a_k)u,\dd(a_k)v)=g_{gH}(u,v)$.
Powstały tensor jest więc $G$-niezmienniczy.
Jego dodatnia określoność wynika z tego, że różniczki działań
są izomorfizmami. W lokalnym przekroju $s:U\to G$ możemy
podstawić $g=s(x)$ do tego samego wzoru;
z gładkości działania i $s$ wynika gładkość metryki na $U$.
\end{proof}

Warunek ze stwierdzenia pozwala także wykryć, ile metryk
jest możliwych. Dla $E(n)/\mathrm O(n)$ zapiszmy iloczyn
skalarny w zerze jako $b(u,v)=u^{\mathsf T}Bv$.
Niezmienniczość wymaga $Q^{\mathsf T}BQ=B$ dla każdego
$Q\in\mathrm O(n)$. Odbicie zmieniające znak tylko $i$-tej
współrzędnej zmienia znak każdego wyrazu $B_{ij}$ z $j\ne i$,
więc wszystkie takie wyrazy są zerowe.
Macierze zamieniające wektory bazy wymuszają z kolei
$B_{ii}=B_{jj}$. Zatem $B=cI$ z $c>0$:
wszystkie niezmiennicze metryki euklidesowe różnią się
jedynie skalą. Dla $n=1$ ostatni wniosek wynika od razu
z wymiaru przestrzeni stycznej.

\section{Zwarte stabilizatory, właściwe działania i ilorazy}
\label{sec:klein-zwarte-stabilizatory}

W przestrzeni jednorodnej $X\cong G/H$ stabilizator $H$
opisuje ruchy pozostawiające wybrany punkt na miejscu.
Zwartość $H$ pozwala uśrednić iloczyn skalarny nad tymi ruchami.
Na przykład stabilizator zera w grupie ruchów euklidesowych
$\R^n\rtimes\mathrm O(n)$ jest zwarty: $\mathrm O(n)$ jest
domkniętym i ograniczonym zbiorem macierzy. Dla pełnej grupy
afinicznej stabilizatorem jest niezwarta $\mathrm{GL}(n,\R)$,
zawierająca wszystkie dylatacje $tI$, $t>0$.
Oba działania obliczymy w
Przykładach~\ref{prz:klein-euklidesowa}
i~\ref{prz:klein-afiniczna}.

Potrzebujemy jednego wyniku zewnętrznego o grupach zwartych:
na każdej zwartej grupie Liego $H$ istnieje \emph{znormalizowana
miara Haara} $\mu$, czyli miara probabilistyczna niezmiennicza
przy mnożeniu z lewej i z prawej strony. Istnienie takiej miary
pochodzi z twierdzenia Haara; zob.
\href{https://doi.org/10.2307/1968346}
{A.\ Haar, \emph{Der Massbegriff in der Theorie der kontinuierlichen
Gruppen}, Ann. Math. 34 (1933), 147--169}.
Poniższy argument uśredniania występuje także u
\href{https://sschleimer.warwick.ac.uk/Teach/2021_3MFDS/References/1983the_geometries_of_three-manifolds.pdf}
{P.\ Scotta, \emph{The Geometries of 3-Manifolds}, \S\,5}.

\begin{stwierdzenie}[Metryka ze zwartego stabilizatora]
\label{prop:klein-zwarty-stabilizator-metryka}
Niech grupa Liego $G$ działa gładko i przechodnio na spójnej
rozmaitości $X$ bez brzegu. Jeśli stabilizator $H=G_{x_0}$ jest
zwarty, to na $X$ istnieje $G$-niezmiennicza metryka
riemannowska.
\end{stwierdzenie}
\begin{proof}
Oznaczmy przez $\rho(h)=\dd(a_h)_{x_0}$ działanie $H$
na $T_{x_0}X$. Wybierzmy dowolny dodatnio określony iloczyn
skalarny $b$ na tej przestrzeni i zdefiniujmy
\[
 b_H(u,v)=\int_H b\bigl(\rho(h)u,\rho(h)v\bigr)\,\dd\mu(h).
\]
Całka jest dobrze określona, bo całkujemy ciągłą funkcję po
zwartej grupie. Jest symetryczna i dwuliniowa. Jeśli $u\ne0$,
to $\rho(h)u\ne0$ dla każdego $h$, więc każdy składnik
$b(\rho(h)u,\rho(h)u)$ jest dodatni.
Ta ciągła funkcja na zwartej $H$ ma dodatnie minimum $m$,
więc $b_H(u,u)\geq m\mu(H)=m>0$.
Niezmienniczość miary przy mnożeniu z prawej daje, dla $k\in H$,
\begin{align*}
 b_H\bigl(\rho(k)u,\rho(k)v\bigr)
 &=\int_H b\bigl(\rho(hk)u,\rho(hk)v\bigr)\,\dd\mu(h)\\
 &=\int_H b\bigl(\rho(\ell)u,\rho(\ell)v\bigr)\,\dd\mu(\ell)
 =b_H(u,v).
\end{align*}
W drugim przejściu podstawiliśmy $\ell=hk$;
całkujemy nadal po całej grupie $H$, a
$(h\mapsto hk)_*\mu=\mu$ jest dokładnie prawostronną
niezmienniczością miary Haara.
Zatem $b_H$ spełnia warunek
Stwierdzenia~\ref{prop:klein-metryka-niezmiennicza}:
przenosząc ten iloczyn skalarny z $x_0$ do punktu
$a\cdot x_0$ przez różniczkę działania $a\in G$, otrzymujemy
metrykę $G$-niezmienniczą. Wybór $a$ nie ma znaczenia,
ponieważ dwa takie elementy różnią się o element $H$.
Gładkość wynika z lokalnych przekrojów submersji $G\to G/H$.
\end{proof}

Zwartość stabilizatora ma też znaczenie topologiczne: zapobiega
temu, by elementy grupy uciekające do nieskończoności przenosiły
punkty pozostające w ustalonym zwartym obszarze na inne punkty
tego samego obszaru.

\begin{definicja}[Działanie właściwe]
\label{def:klein-dzialanie-wlasciwe}
\index{dzialanie wlasciwe@Działanie właściwe}
Niech grupa Liego $G$ działa na rozmaitości $X$.
Działanie jest \emph{właściwe}, gdy odwzorowanie
\[
 A:G\times X\longrightarrow X\times X,
 \qquad A(a,x)=(x,a\cdot x),
\]
jest właściwe, to znaczy przeciwobraz każdego zbioru zwartego
jest zwarty. Jeżeli $G$ jest grupą dyskretną, równoważnie
dla każdego zwartego $K\subseteq X$ zbiór
\[
 \{a\in G:aK\cap K\ne\varnothing\}
\]
jest skończony. W tej drugiej postaci warunek mówi, że tylko
skończenie wiele ruchów może sprowadzić zwarty obszar z powrotem
do siebie.
\end{definicja}

Istotnie, przy właściwym działaniu wskazany zbiór jest obrazem
zwartego $A^{-1}(K\times K)$ przez rzut na dyskretną grupę,
a więc jest skończony. Odwrotnie, jeśli jest skończony dla
każdego $K$, przeciwobraz zwartego $C\subset X\times X$ leży
w skończonej sumie zbiorów $\{a\}\times K$, gdzie $K$ zawiera
oba rzuty $C$; jest w niej domknięty, zatem zwarty.

\begin{stwierdzenie}[Właściwość działania na $G/H$]
\label{prop:klein-zwarte-h-wlasciwe}
Niech grupa Liego $G$ działa gładko i przechodnio na spójnej
rozmaitości $X$ bez brzegu, a $H=G_{x_0}$.
Działanie $G$ na $X$ jest właściwe wtedy i tylko wtedy,
gdy $H$ jest zwarty.
\end{stwierdzenie}
\begin{proof}
Jeśli działanie jest właściwe, to
$A^{-1}(\{(x_0,x_0)\})=H\times\{x_0\}$ jest zwarty,
więc $H$ jest zwarty.

Załóżmy odwrotnie, że $H$ jest zwarty. Na mocy
Twierdzenia~\ref{thm:klein-orbita-iloraz} możemy utożsamić
$X$ z $G/H$. Niech $C\subseteq X\times X$ będzie zwarty,
a $K_1,K_2\subseteq X$ jego rzutami na oba czynniki.
Submersja $q:G\to G/H$ ma lokalne gładkie przekroje.
Z lokalnej zwartości $X$ możemy wybrać skończenie wiele
otwartych zbiorów $U_i$ pokrywających $K_1$ oraz $V_j$
pokrywających $K_2$, tak aby ich domknięcia były zwarte
i mieściły się w obszarach odpowiednich lokalnych przekrojów
$s_i$ oraz $t_j$ odwzorowania $q$.
Jeśli $(a,x)\in A^{-1}(C)$, to dla pewnych $i,j$ mamy
$x\in U_i$ i $y=a\cdot x\in V_j$. Ponieważ
$s_i(x)H=x$ i $t_j(y)H=y$, zachodzi
\[
 h:=t_j(y)^{-1}a s_i(x)\in H,
 \qquad
 a=t_j(y)h s_i(x)^{-1}.
\]
Zatem wszystkie możliwe $a$ leżą w skończonej sumie
zbiorów zwartych
\[
 L=\bigcup_{i,j}t_j(\overline V_j)\,
 H\,s_i(\overline U_i)^{-1}\subseteq G.
\]
Przeciwobraz $A^{-1}(C)$ jest domkniętym podzbiorem
zwartego $L\times K_1$, a więc jest zwarty.
\end{proof}

Ważne jest tu założenie przechodniości. Działanie $\R$ na
torusie przez nawinięcie niewymierne, opisane przed
Definicją~\ref{def:klein-geometria}, ma trywialne, a więc
zwarte stabilizatory, ale nie jest właściwe: jego orbita jest
gęsta i nie jest domknięta. Przy właściwym działaniu
odwzorowanie orbity $a\mapsto a\cdot x$ jest właściwe
(ograniczamy $A$ do $G\times\{x\}$), a więc ma domknięty obraz:
jeśli $a_j\cdot x\to y$, właściwość daje podciąg $a_j\to a$
po ograniczeniu do zwartego otoczenia $y$, skąd $y=a\cdot x$.

\begin{stwierdzenie}[Zupełność metryki jednorodnej]
\label{prop:klein-jednorodna-zupelna}
Jeśli spójna rozmaitość riemannowska $(X,g)$ jest jednorodna,
czyli jej grupa izometrii działa przechodnio, to $(X,g)$ jest
geodezyjnie i metrycznie zupełna.
\end{stwierdzenie}
\begin{proof}
Wybierzmy $x_0\in X$. Lokalna teoria równania geodezyjnej
daje $\varepsilon>0$ takie, że $\exp_{x_0}$ jest określone
na $\{v\in T_{x_0}X:|v|_g<\varepsilon\}$.
Ustalmy $0<\delta<\varepsilon$. Dla każdego $x\in X$
istnieje izometria $a$ z $a(x_0)=x$. Izometria przenosi
geodezyjne na geodezyjne i zachowuje długość wektorów,
więc $\exp_x$ jest określone na kuli $|v|_g<\varepsilon$.
Każdą geodezyjną o prędkości $1$ można zatem przedłużyć
od dowolnego jej punktu co najmniej o czas $\delta$.
Gdyby jej maksymalny przedział miał skończony prawy koniec $b$,
wybralibyśmy chwilę $t_0>b-\delta$; przedłużenie od
$\gamma(t_0)$ ma postać
$s\mapsto\exp_{\gamma(t_0)}(s\dot\gamma(t_0))$ dla $|s|<\varepsilon$.
Istnieje w szczególności do $s=\delta$ i na wspólnym
przedziale zgadza się z $s\mapsto\gamma(t_0+s)$
przez jednoznaczność równania geodezyjnej.
Sięga więc poza $b$, co jest sprzecznością.
Podobnie postępujemy na lewym końcu. Dla innej niezerowej
stałej prędkości przeskalowujemy parametr, a geodezyjne
o prędkości zerowej są stałe. Zupełność metryczna wynika
z Twierdzenia Hopfa--Rinowa~\ref{thm:hopf-rinow-geodezyjne}.
\end{proof}

Pozostaje odróżnić sam model $X$ od rozmaitości otrzymanej
przez utożsamienie punktów modelu. Wymagamy, by grupa
utożsamień była dyskretna w topologii grupy Liego i działała
bez punktów stałych.

\begin{definicja}[Podgrupa dyskretna i działanie swobodne]
\label{def:klein-dyskretna-swobodna}
\index{podgrupa dyskretna}
\index{dzialanie swobodne@Działanie swobodne}
Podgrupa $\Gamma<G$ jest \emph{dyskretna}, jeśli jedność
$e\in G$ ma otoczenie $V$ spełniające $V\cap\Gamma=\{e\}$.
Działanie $\Gamma$ na $X$ jest \emph{swobodne} (także: wolne), jeśli
$\gamma\cdot x=x$ implikuje $\gamma=e$ dla każdych
$x\in X$ i $\gamma\in\Gamma$.
Na przykład $\Z^n<\R^n\rtimes\mathrm O(n)$ jest dyskretną
podgrupą translacji i działa swobodnie na $\R^n$.
\end{definicja}

Podgrupa dyskretna grupy Liego jest też domknięta.
Gdyby różne $\gamma_j\in\Gamma$ zbiegały do punktu grupy $G$,
to $\gamma_j^{-1}\gamma_{j+1}\to e$, mimo że są to niejedynkowe
elementy $\Gamma$. Przeczyłoby to definicji dyskretności.

\begin{twierdzenie}[Iloraz przez swobodną grupę dyskretną]
\label{prop:klein-dyskretny-iloraz}
Niech $G$ działa gładko i przechodnio na spójnej
rozmaitości $X$ bez brzegu, ze zwartym stabilizatorem.
Niech $g$ będzie metryką $G$-niezmienniczą, a
$\Gamma<G$ dyskretną podgrupą działającą swobodnie na $X$.
Wtedy $M=\Gamma\backslash X$ jest gładką rozmaitością,
rzutowanie $\pi:X\to M$ jest nakryciem i istnieje
jednoznaczna metryka $g_M$, dla której $\pi$ jest lokalną
izometrią. Rozmaitości $(X,g)$ i $(M,g_M)$ są zupełne.
Jeśli $X$ jest jednospójna, $\pi$ jest nakryciem uniwersalnym,
a po wyborze punktów bazowych $\pi_1(M)\cong\Gamma$.
\end{twierdzenie}
\begin{proof}
Z domkniętości $\Gamma$ i
Stwierdzenia~\ref{prop:klein-zwarte-h-wlasciwe}
wynika, że działanie $\Gamma$ jest właściwe: przeciwobraz
zwartego $C\subset X\times X$ przez
$\Gamma\times X\to X\times X$ jest domkniętym podzbiorem
zwartego przeciwobrazu przez $G\times X\to X\times X$.

Ustalmy $x\in X$. Pokażemy, że istnieje otoczenie $U$ punktu $x$
takie, iż $\gamma U\cap U=\varnothing$ dla każdego
$\gamma\ne e$. Gdyby nie istniało, moglibyśmy wybrać
malejące do $x$ otoczenia $U_j$ o wspólnym zwartym domknięciu,
punkty $u_j,v_j\in U_j$ i $\gamma_j\ne e$ z
$\gamma_j u_j=v_j$. Wtedy $u_j,v_j\to x$.
Właściwość działania grupy dyskretnej oznacza, że tylko
skończenie wiele $\gamma_j$ może przenieść to zwarte
domknięcie na siebie. Na podciągu mamy więc stałe
$\gamma_j=\gamma\ne e$, a przejście do granicy daje
$\gamma x=x$, wbrew swobodzie działania.
Zbiory $\gamma U$ są parami rozłączne i
\[
 \pi^{-1}(\pi(U))=\coprod_{\gamma\in\Gamma}\gamma U.
\]
Każde ograniczenie $\pi|_{\gamma U}$ jest homeomorfizmem
na $\pi(U)$, więc $\pi$ jest nakryciem.

Iloraz jest Hausdorffa. Jeśli dwie różne orbity $\Gamma x$,
$\Gamma y$ nie miałyby rozłącznych otoczeń w ilorazie,
wybralibyśmy malejące otoczenia $U_j\ni x$, $V_j\ni y$
o zwartych domknięciach zawartych w ustalonym zbiorze zwartym
oraz $u_j\in U_j$, $v_j\in V_j$ i $\gamma_j\in\Gamma$
z $\gamma_j u_j=v_j$. Właściwość działania znów pozwala
wybrać stałe $\gamma_j=\gamma$ na podciągu, skąd
$y=\gamma x$, sprzeczność. Rzutowanie $\pi$ jest otwarte,
bo $\pi^{-1}(\pi(U))=\bigcup_{\gamma\in\Gamma}\gamma U$
jest otwarte dla każdego otwartego $U\subset X$; dlatego
obrazy przeliczalnej bazy $X$ tworzą przeliczalną bazę
ilorazu. Wybierając $U$ wewnątrz obszaru mapy $X$,
przenosimy tę mapę przez $\pi|_U$ na $\pi(U)$;
przejścia między otrzymanymi mapami są lokalnie ograniczeniami
działania elementów $\Gamma$ i są gładkie.

Dla $[x]=\pi(x)$ definiujemy
\[
 (g_M)_{[x]}\bigl(\dd\pi_x u,\dd\pi_x v\bigr)
 :=g_x(u,v).
\]
To nie zależy od przedstawiciela $x$: jeśli $x'=\gamma x$,
to $\pi\circ\gamma=\pi$, a $\gamma$ jest izometrią $g$.
Różniczkując tę równość, dostajemy
$\dd\pi_{x'}\circ\dd\gamma_x=\dd\pi_x$.
Zatem te same wektory w $T_{[x]}M$ mają w drugim
przedstawicielu podniesienia $\dd\gamma_xu,\dd\gamma_xv$,
a ich iloczyn skalarny wynosi
$g_{x'}(\dd\gamma_xu,\dd\gamma_xv)=g_x(u,v)$.
W lokalnej mapie nakrycia wzór jest przeniesieniem
gładkiej, dodatnio określonej metryki, więc określa
metrykę gładką; warunek lokalnej izometrii wymusza jej wartości,
co dowodzi jednoznaczności.
Rzutowanie jest więc lokalną izometrią.
Zupełność $X$ daje
Stwierdzenie~\ref{prop:klein-jednorodna-zupelna}.
Każdą geodezyjną w $M$ podnosimy do geodezyjnej w $X$;
przedłużenie podniesionej geodezyjnej rzutuje się na
przedłużenie w $M$. Zatem $M$ jest geodezyjnie zupełna,
a z Twierdzenia~\ref{thm:hopf-rinow-geodezyjne} również
metrycznie zupełna. Jeśli $X$ jest jednospójna, nakrycie
$\pi$ jest uniwersalne. Aby zidentyfikować grupę podstawową,
ustalmy $x_0\in X$. Pętla w $M$ oparta w $[x_0]$ podnosi się
od $x_0$ do drogi kończącej się w dokładnie jednym punkcie
$\gamma x_0$, $\gamma\in\Gamma$. Przypisanie pętli tego
elementu $\gamma$ jest homomorfizmem: podniesienie drugiej
pętli od $\gamma x_0$ jest obrazem jej podniesienia od $x_0$
przez działanie $\gamma$. Każdy $\gamma$ jest osiągany, bo
$X$ jest łukowo spójna i istnieje droga od $x_0$ do
$\gamma x_0$. Jądro składa się z pętli o zamkniętym
podniesieniu, które są zerowe, ponieważ $X$ jest jednospójna
i homotopie podnoszą się przez nakrycie. Jest to teoria
podnoszenia z Rozdziału~\ref{ch:topologia-algebraiczna-podstawy}.
\end{proof}

Przy translacjach $\Gamma=\Z^n$ na $\R^n$ dostajemy płaski
torus $T^n=\Z^n\backslash\R^n$. Warunek swobody jest istotny:
odbicie $\R^n$ w hiperpłaszczyźnie ma punkty stałe, więc jego
iloraz nie jest nakryty przez $\R^n$ w punktach obrazu tej
hiperpłaszczyzny. Samo $G$ nie musi natomiast działać na
$\Gamma\backslash X$: jego elementy mogą nie zachowywać
orbit $\Gamma$.

Jeśli stabilizatory nie są trywialne, lokalny iloraz zachowuje
informację o symetriach ustalających punkt. Tę informację zapisuje
struktura orbifoldu. Jest ona potrzebna w opisie baz rozwłóknień
Seiferta w Rozdziale~\ref{ch:geometrizacja-thurstona}.

\begin{definicja}[Orbifold gładki]
\label{def:klein-orbifold}
\index{orbifold}
\index{orbifold!grupa lokalna}
\index{orbifold!przestrzeń podkładowa}
\emph{Efektywnym orbifoldem gładkim wymiaru $n$} nazywamy
przestrzeń topologiczną $|\mathcal O|$, która jest Hausdorffa i ma
przeliczalną bazę, wyposażoną w zgodny atlas \emph{map
orbifoldowych}. Mapa wokół $y\in|\mathcal O|$ składa się z otwartej
kuli $B\subset\R^n$ o środku $0$, skończonej grupy
$H\leq\mathrm O(n)$ i homeomorfizmu
\[
 \varphi:B/H\longrightarrow U\subset|\mathcal O|,
 \qquad \varphi([0])=y.
\]
Symbol $[v]_H$ oznacza orbitę punktu $v$ pod działaniem $H$;
na $B/H$ bierzemy topologię ilorazową.
Obrazy takich map są otwarte i pokrywają $|\mathcal O|$.
Wyjaśnijmy dokładnie, co znaczy \emph{zgodność} map.
Niech $(C,K,\psi)$ i $(B,H,\varphi)$ będą dwiema mapami,
o obrazach odpowiednio $V\subset U$. \emph{Zanurzeniem pierwszej
mapy w drugiej} jest para złożona z injektywnego homomorfizmu
$\iota:K\hookrightarrow H$ i gładkiego zanurzenia
$f:C\hookrightarrow B$ spełniająca dla każdego $v\in C$
i $k\in K$ dwa warunki:
\[
 f(kv)=\iota(k)f(v),
 \qquad
 \varphi\bigl([f(v)]_H\bigr)=\psi\bigl([v]_K\bigr).
\]
Pierwszy warunek mówi, że $f$ zachowuje działanie grup po
identyfikacji $K$ z podgrupą $H$. Drugi mówi, że po przejściu
do ilorazów $f$ daje zwykłą inkluzję $V\hookrightarrow U$.
Ponieważ obie kule mają wymiar $n$, zanurzenie $f$ jest
dyfeomorfizmem na otwarty obraz.

Atlas jest \emph{zgodny}, gdy dla każdych dwóch jego map
o obrazach $U_1,U_2$ i każdego $y\in U_1\cap U_2$
istnieje trzecia mapa $(C,K,\psi)$, której obraz $V$ spełnia
$y\in V\subset U_1\cap U_2$ i która zanurza się w obie
początkowe mapy. Wybieramy więc mniejszy model całego
otoczenia $y$, a nie tylko odwzorowanie między przestrzeniami
podkładowymi. Dwa atlasy określają tę samą strukturę, gdy
ich połączenie jest zgodne. Zmiana podniesienia $f$ przez
element $h\in H$ zmienia $\iota$ przez sprzężenie
$\iota(k)\mapsto h\iota(k)h^{-1}$; opisuje tę samą inkluzję
na ilorazach.

W mapie wyśrodkowanej w $y$ grupa $H$ jest \emph{grupą lokalną}
tego punktu. Jeśli $z\in U$ odpowiada klasie punktu $v\in B$,
jego grupą lokalną jest stabilizator $H_v$; inny wybór punktu
w tej samej klasie sprzęga stabilizator w $H$. Warunek zgodności
zapewnia, że grupa lokalna nie zależy, z dokładnością do
izomorfizmu, od wybranej mapy. Jej rząd mierzy miejscową
symetrię. Sama przestrzeń $|\mathcal O|$ to \emph{przestrzeń
podkładowa} orbifoldu. Jeśli wszystkie grupy lokalne są
trywialne, otrzymujemy zwykły atlas gładkiej rozmaitości.
\end{definicja}
Jest to wersja efektywna definicji z
\href{https://library.slmath.org/nonmsri/gt3m/PDF/13.pdf}
{W.\ P.\ Thurstona, \emph{The Geometry and Topology of
Three-Manifolds}, rozdział 13, \S\,13.2}.

\begin{przyklad}[Rozmaitość i punkty osobliwe orbifoldu]
\label{prz:klein-orbifoldy}
Każda gładka rozmaitość jest orbifoldem z grupami lokalnymi
$\{e\}$. Odwrotna implikacja nie zachodzi. Dla $m\geq2$ niech
$D=\{z\in\C:|z|<1\}$ i niech $C_m$ działa na nim przez obroty
$z\mapsto e^{2\pi i/m}z$. Iloraz $D/C_m$ ma w obrazie środka
grupę lokalną $C_m$, a poza nim grupy trywialne:
jeśli $\zeta^kz=z$ i $z\ne0$, to $\zeta^k=1$.
Odwzorowanie $z\mapsto z^m$ jest stałe na orbitach,
a równość $z^m=w^m$ oznacza $w=\zeta^kz$ dla pewnego $k$.
Każdy punkt dysku ma pierwiastek stopnia $m$ w dysku.
Otrzymujemy więc ciągłą bijekcję
$[z]\mapsto z^m$ z $|D/C_m|$ na dysk.
Jest ona homeomorfizmem: poza zerem wybieramy lokalną
gałąź pierwiastka, a w zerze ciągłość odwrotności
wynika z $|z|=|z^m|^{1/m}$.
Ta identyfikacja jest topologiczna,
ale nie usuwa grupy lokalnej: jako orbifold iloraz nie jest zwykłą
rozmaitością. Standardowa metryka ilorazowa ma przy środku kąt
stożkowy $2\pi/m$. Rzeczywiście, reprezentantów orbit
wybieramy w sektorze $0\leq\theta\leq2\pi/m$, sklejając
jego boki promieniowe. Metryka poza wierzchołkiem ma postać
$\dd r^2+r^2\dd\theta^2$, więc okrąg promienia $r$ ma długość
$\int_0^{2\pi/m}r\,\dd\theta=2\pi r/m$.
Iloraz obwodu przez promień daje podany kąt.

Zobaczmy warunek zgodności na tych mapach. Niech
$q:D\to D/C_m$ będzie rzutem, a $a\in D\setminus\{0\}$.
Wybierzmy $\varepsilon>0$ tak małe, by dyski
$\zeta^k(a+B_\varepsilon(0))$, gdzie
$\zeta=e^{2\pi i/m}$ i $0\leq k<m$, były parami rozłączne.
Wystarczy
$\varepsilon<\frac12\min_{1\leq k<m}|a-\zeta^ka|$
i $\varepsilon<1-|a|$: pierwsze minimum jest dodatnie
i ogranicza od dołu odległości wszystkich różnych środków.
Mapa wokół $q(a)$ ma wtedy postać
$(B_\varepsilon(0),\{e\},\psi)$, gdzie
$\psi(v)=q(a+v)$. Zanurza się ona w mapę
$(D,C_m,\operatorname{id}_{D/C_m})$ przez $f(v)=a+v$
i jedyny homomorfizm $\{e\}\hookrightarrow C_m$:
równość $q(f(v))=\psi(v)$ jest dokładnie drugim
warunkiem z definicji.
Wokół $q(0)$ bierzemy mniejszy dysk o środku $0$ z tą
samą grupą $C_m$; tu zanurzenie map jest zwykłą
inkluzją dysków i identycznością grup.

Odbicie $z\mapsto\bar z$ daje iloraz $D$ z \emph{krawędzią
zwierciadlaną}: jej punkty mają stabilizator $C_2$, a przestrzeń
podkładowa jest półdyskiem. W wyższym wymiarze iloraz
$\R^3/\{\pm I\}$ ma przy obrazie zera przekrój $\mathbb{RP}^2$.
Przebite otoczenie tego punktu ma więc grupę podstawową $\Z/2$,
bo ma postać $(0,r)\times\mathbb{RP}^2$ i deformacyjnie
retrahuje się na $\mathbb{RP}^2$.
Włączenia mniejszych takich otoczeń w większe indukują
izomorfizmy tej grupy. Gdyby istniała mapa rozmaitości
w obrazie zera, między dwoma takimi otoczeniami moglibyśmy
wstawić otoczenie homeomorficzne z kulą.
Po usunięciu środka włączenie faktoryzowałoby się przez
przebitą kulę w $\R^3$, która jest jednospójna,
i indukowałoby homomorfizm zerowy zamiast izomorfizmu.
Tutaj nawet
przestrzeń podkładowa nie jest rozmaitością w obrazie zera.
\end{przyklad}

\begin{figure}[htbp]
\centering
\begin{tikzpicture}[>=Stealth,scale=.94]
 \begin{scope}
  \fill[manifoldblue!14] (0,0)--(1.35,0)
   arc[start angle=0,end angle=120,radius=1.35]--cycle;
  \draw[manifoldblue!45] (0,0) circle(1.35);
  \foreach \angle in {0,120,240}
   {\draw[manifoldblue!30,dashed] (0,0)--(\angle:1.35);}
  \draw[->,thick,manifoldred] (0,0)--(1.28,0);
  \draw[->,thick,manifoldred] (0,0)--(120:1.28);
  \draw[manifoldgold!85!black] (.43,0)
   arc[start angle=0,end angle=120,radius=.43];
  \node at (.73,.57) {$2\pi/3$};
  \fill[manifoldblue!80!black] (0,0) circle(2pt);
  \node[below left=3pt] at (0,0) {$[0]$};
  \node[align=center,font=\small] at (0,-1.95)
   {$D/C_3$: boki sektora sklejamy;\\wierzchołek ma grupę $C_3$};
 \end{scope}
 \begin{scope}[xshift=6.4cm]
  \fill[manifoldblue!14] (-1.35,0)
   arc[start angle=180,end angle=0,radius=1.35]--cycle;
  \draw[manifoldblue!45] (0,0) circle(1.35);
  \draw[very thick,manifoldred] (-1.35,0)--(1.35,0);
  \fill[manifoldblue] (.6,.7) circle(2pt) (.6,-.7) circle(2pt);
  \draw[->,manifoldgold!85!black] (.6,-.6)--(.6,.6);
  \node[right=4pt] at (.6,.7) {$z$};
  \node[right=4pt] at (.6,-.7) {$\bar z$};
  \node[above=3pt,font=\small] at (-.35,0) {$C_2$};
  \node[align=center,font=\small] at (0,-1.95)
   {$D/\langle z\mapsto\bar z\rangle$: półdysk;\\średnica jest zwierciadłem};
 \end{scope}
\end{tikzpicture}
\caption{Dwa lokalne modele orbifoldowe. Po lewej pozostaje
jedna trzecia kąta pełnego; po prawej odbicie utożsamia
dolną połowę dysku z górną i ustala średnicę.
Okręgi wyznaczają zasięg rysunku; brzeg kołowy nie należy do modelu.}
\label{fig:klein-orbifold-sektor-zwierciadlo}
\end{figure}

\begin{stwierdzenie}[Iloraz z punktami stałymi]
\label{prop:klein-iloraz-orbifoldowy}
Niech $G$ działa gładko i przechodnio na spójnej rozmaitości $X$
ze zwartym stabilizatorem, a $\Gamma<G$ będzie dyskretna i działa
skutecznie na $X$. Wtedy $\Gamma\backslash X$ ma naturalną
strukturę efektywnego orbifoldu. Dla $x\in X$ jego grupa lokalna
jest izomorficzna ze stabilizatorem $\Gamma_x$.
\end{stwierdzenie}
\begin{proof}
Z własności działania $G$ na $X$ wynika właściwość działania
$\Gamma$, jak w dowodzie
Twierdzenia~\ref{prop:klein-dyskretny-iloraz}. Przeciwobraz
punktu $(x,x)$ przez odwzorowanie działania jest zwarty, więc
$\Gamma_x$ jest zwarta; jako podgrupa dyskretna jest skończona.
Właściwość działania daje przestrzeń orbit Hausdorffa,
a obrazy przeliczalnej bazy $X$ przez otwarty rzut na iloraz
tworzą jej przeliczalną bazę, jak w przywołanym dowodzie.
Wybierzmy niezmienniczą metrykę ze
Stwierdzenia~\ref{prop:klein-zwarty-stabilizator-metryka}.
Uzasadnijmy, że orbitę $x$ można odizolować od jej pozostałych
punktów. Wybierzmy $\varepsilon>0$ tak, aby domknięta kula
$\overline B_\varepsilon(x)$ była zwarta.
Właściwość działania daje skończony zbiór
$S=\{\gamma:\gamma\overline B_\varepsilon(x)
\cap\overline B_\varepsilon(x)\ne\varnothing\}$.
Każdy $\gamma$ z $d(x,\gamma x)\leq\varepsilon$ należy do $S$,
bo $x$ leży w kuli, a $\gamma x$ w jej obrazie i w samej kuli.
Jeśli $S\setminus\Gamma_x$ jest niepusty, skończenie wiele
odległości $d(x,\gamma x)$ dla $\gamma$ z tego zbioru
ma dodatnie minimum. Bierzemy $2r$ mniejsze od
$\varepsilon$ i od wszystkich tych odległości;
jeśli zbiór jest pusty, wystarcza $2r<\varepsilon$.
Wtedy $d(x,\gamma x)>2r$ dla każdego $\gamma\notin\Gamma_x$.
Zmniejszamy $r$ tak, by
odwzorowanie wykładnicze było dyfeomorfizmem na kuli
$B_r(0)\subset T_xX$. Wówczas geodezyjna kula $B_r(x)$ jest
zachowywana przez $\Gamma_x$, a jej obrazy przez elementy spoza
$\Gamma_x$ są z nią rozłączne.
Istotnie, punkt w przecięciu $B_r(x)\cap\gamma B_r(x)$
dawałby z nierówności trójkąta $d(x,\gamma x)<2r$,
ponieważ $\gamma$ jest izometrią.
Rzut na iloraz identyfikuje zatem
otoczenie orbity $x$ z $B_r(x)/\Gamma_x$.
Ponieważ dla $h\in\Gamma_x$ zachodzi
$h\exp_x(v)=\exp_x(\dd h_xv)$, otrzymujemy model
$B_r(0)/\Gamma_x$ z ortogonalnym działaniem. To działanie jest
skuteczne: jeśli $\dd h_x=\id$, równość z eksponentą daje
$h=\id$ na otoczeniu $x$. Zbiór punktów $y$ spełniających
$h(y)=y$ i $\dd h_y=\id$ jest otwarty dzięki tej samej
równości w otoczeniach normalnych oraz domknięty przez
ciągłość. Spójność $X$ wymusza $h=\id_X$;
skuteczność działania $\Gamma$ daje wtedy $h=e$.
Sprawdźmy zgodność tak otrzymanych map. Jeśli ich obrazy
przecinają się w orbicie $[z]$, wybieramy przedstawicieli
$z_1\in B_r(x)$ i $z_2\in B_s(y)$ tej orbity oraz
$\gamma\in\Gamma$ takie, że $z_2=\gamma z_1$.
Po przeniesieniu drugiej kuli przez $\gamma^{-1}$ obie
zawierają $z_1$. Stabilizator $\Gamma_{z_1}$ jest podgrupą
zarówno $\Gamma_x$, jak i $\gamma^{-1}\Gamma_y\gamma$:
element ustalający $z_1$ nie może przenieść żadnej z tych
kul na rozłączny z nią obraz. Wybieramy małą kulę
geodezyjną $W$ wokół $z_1$, zawartą w przecięciu obu kul
i rozłączną ze swoimi obrazami przez elementy spoza
$\Gamma_{z_1}$. Używając współrzędnych wykładniczych w
$z_1$, otrzymujemy z $W/\Gamma_{z_1}$ wspólną mapę
orbifoldową, która zanurza
się w obie wyjściowe mapy przez inkluzję do pierwszej kuli
oraz przez $\gamma$ po inkluzji do drugiej. Towarzyszące
homomorfizmy grup to odpowiednio inkluzja
$\Gamma_{z_1}\hookrightarrow\Gamma_x$ i odwzorowanie
$h\mapsto\gamma h\gamma^{-1}$ z $\Gamma_{z_1}$ do
$\Gamma_y$.
Ekwiwariantność wynika z podanych homomorfizmów
stabilizatorów. Po rzutowaniu do $\Gamma\backslash X$
obydwa zanurzenia przypisują punktowi $w\in W$ tę samą
orbitę $[w]$, więc spełniają także drugie równanie
z definicji.
\end{proof}

\section{Struktury lokalnie modelowane na przestrzeni jednorodnej}
\label{sec:thur-dop-struktury-gx}
\index{struktura (G,X)}
Geometria Kleina opisuje pojedynczą przestrzeń $X$ z grupą
symetrii $G$. Natomiast rozmaitość może tylko \emph{lokalnie}
wyglądać jak $X$: jej topologia nie musi dopuszczać
jednej globalnej mapy do tego modelu. Formalizujemy to
analogicznie do atlasu rozmaitości gładkiej, lecz ograniczamy
przekształcenia przejścia.

\begin{definicja}[Atlas struktury $(G,X)$]
\label{def:thur-dop-gx-atlas}
Niech $X$ będzie spójną gładką rozmaitością, na której grupa Liego
$G$ działa skutecznie przez dyfeomorfizmy. \emph{Strukturą $(G,X)$}
na rozmaitości $M$ bez brzegu, wymiaru $\dim X$, nazywamy atlas odwzorowań
$\varphi_i:U_i\to V_i\subset X$, będących dyfeomorfizmami
na otwarte podzbiory modelu, taki że na każdej spójnej składowej
przecięcia $U_i\cap U_j$ odwzorowanie przejścia
$\varphi_j\circ\varphi_i^{-1}$ jest ograniczeniem
pewnego elementu $g_{ji}\in G$. Zwykle utożsamiamy atlasy
zgodne w tym sensie.
\end{definicja}

Właściwą klasą lokalnych przekształceń jest pseudogrupa
utworzona z ograniczeń elementów $G$ i sklejania takich
ograniczeń. Zwykły atlas gładki odpowiada znacznie większej
pseudogrupie wszystkich lokalnych dyfeomorfizmów;
w atlasie geometrycznym przejścia muszą zachowywać
wybraną strukturę modelu.
Jeśli $X$ ma $G$-niezmienniczą metrykę $g_X$, określamy
na $U_i$ tensor $g_i=\varphi_i^*g_X$.
Na wspólnej dziedzinie $\varphi_j=g_{ji}\circ\varphi_i$,
więc
$g_j=\varphi_i^*(g_{ji}^*g_X)=\varphi_i^*g_X=g_i$.
Tak otrzymujemy metrykę na całej $M$.
Istnienie atlasu geometrycznego nie przesądza jeszcze,
czy ta metryka jest zupełna.

\begin{twierdzenie}[Rozwinięcie i reprezentacja holonomii]
\label{thm:thur-dop-dev-hol}
Załóżmy, że $X$ jest spójna i jednospójna, a $G$ działa
skutecznie przez izometrie pewnej metryki riemannowskiej
na $X$. Dla spójnej rozmaitości $M$ ze strukturą $(G,X)$
istnieją lokalny dyfeomorfizm
$\mathrm{dev}:\widetilde M\to X$ i homomorfizm
$\rho:\pi_1(M)\to G$, spełniające
\begin{equation}\label{eq:thur-dop-dev-equiv}
 \mathrm{dev}(\gamma\cdot\widetilde p)
 =\rho(\gamma)\cdot\mathrm{dev}(\widetilde p).
\end{equation}
Po wyborze punktów bazowych i identyfikacji przekształceń
nakrycia para $(\mathrm{dev},\rho)$ jest wyznaczona z dokładnością
do jednoczesnego złożenia $\mathrm{dev}$ z elementem $G$
i sprzężenia $\rho$ tym elementem.
\end{twierdzenie}
\begin{proof}
Wybierzmy punkt bazowy i jedną mapę atlasu w jego otoczeniu.
Najpierw zauważmy, że dwie izometrie $X$ zgodne na niepustym
zbiorze otwartym są równe. Ich złożenie odwrotne ustala tam
punkt i jego całą przestrzeń styczną. Izometria spełnia
$F(\exp_p v)=\exp_{F(p)}(\dd F_pv)$, więc te same dane
ustalają ją w otoczeniu normalnym. Zbiór punktów, w których
ustala punkt i różniczkę, jest zatem otwarty; jest też domknięty
przez ciągłość. Spójność $X$ daje równość na całym $X$.
W szczególności element grupy odpowiadający przejściu
jest jednoznaczny na każdej składowej przecięcia.

Przedłużamy początkową mapę wzdłuż drogi, pokrywając jej
zwarty obraz skończenie wieloma dziedzinami map. Jeśli dotychczasowa
mapa to $h_i\circ\varphi_i$, a
$\varphi_j=g_{ji}\circ\varphi_i$ na otoczeniu przejścia,
następną wybieramy jako
$(h_i\circ g_{ji}^{-1})\circ\varphi_j$.
Na otoczeniu punktu przejścia mamy dokładnie
$h_i\circ g_{ji}^{-1}\circ\varphi_j
=h_i\circ g_{ji}^{-1}\circ g_{ji}\circ\varphi_i
=h_i\circ\varphi_i$.
Wstawienie dodatkowej mapy
nie zmienia wyniku, ponieważ na wspólnym otoczeniu
$g_{ki}=g_{kj}\circ g_{ji}$; równość lokalna jest równością
elementów $G$ na mocy poprzedniego akapitu.

Dla homotopii dróg z ustalonymi końcami wybieramy dostatecznie
drobną siatkę w kwadracie parametrów, aby obraz każdego
prostokąta leżał w jednej dziedzinie mapy. Istnieje taka
siatka dzięki lematowi Lebesgue'a. Przesunięcie drogi
z dwóch boków małego prostokąta na dwa pozostałe nie zmienia
przedłużenia: w jednej mapie oba obliczenia są tym samym
odwzorowaniem. Przechodząc przez skończenie wiele prostokątów,
otrzymujemy niezależność od homotopii. Model nakrycia uniwersalnego
przez klasy dróg daje więc $\mathrm{dev}$; lokalnie jest ono
złożeniem mapy atlasu z elementem $G$, czyli lokalnym dyfeomorfizmem.

Zastosowanie tego przepisu do pętli opartej w punkcie
bazowym zmienia początkową mapę na
$h_\gamma\circ\varphi_0$ dla jednoznacznego $h_\gamma\in G$.
Kładziemy $\rho([\gamma])=h_\gamma$.
Sprawdźmy kolejność iloczynu. Zgodnie z Rozdziałem~7
droga $\gamma*\alpha$ przebiega najpierw $\gamma$, potem $\alpha$.
W modelu nakrycia przez klasy dróg działanie lewe ma postać
$T_{[\gamma]}([\alpha])=[\gamma*\alpha]$.
Po przebyciu $\gamma$ zaczynamy przedłużanie wzdłuż $\alpha$
od mapy $h_\gamma\circ\varphi_0$, toteż całe rozwinięcie
wzdłuż $\alpha$ zostaje złożone z lewej z $h_\gamma$.
Zatem
$\mathrm{dev}(T_{[\gamma]}[\alpha])
=h_\gamma\cdot\mathrm{dev}([\alpha])$,
co jest równością~\eqref{eq:thur-dop-dev-equiv}.
Dla dwóch pętli daje to
$h_{\gamma*\delta}=h_\gamma h_\delta$, a więc
$\rho([\gamma][\delta])=\rho([\gamma])\rho([\delta])$.
Inny wybór
początkowej mapy na $a\circ\varphi_0$ daje
$\mathrm{dev}'=a\circ\mathrm{dev}$ oraz
$\rho'(\gamma)=a\circ\rho(\gamma)\circ a^{-1}$.
\end{proof}

\begin{uwaga}[Dwie różne holonomie]
\label{uw:thur-dop-dwie-holonomie}
\index{holonomia!struktury geometrycznej}
Reprezentacja $\rho:\pi_1(M)\to G$ to
\emph{holonomia struktury $(G,X)$}. Nie należy jej
utożsamiać z grupą holonomii koneksji Leviego-Civity
z Definicji~\ref{def:grupa-holonomii}. Na płaskim torusie
$T^n=\R^n/\Z^n$ pierwsza holonomia wysyła element
$m\in\Z^n$ na przesunięcie $x\mapsto x+m$, więc jest
nietrywialna; druga jest trywialna, bo pola współrzędne
schodzą do globalnej ramki równoległej.
\end{uwaga}

\begin{figure}[htbp]
\centering
\begin{tikzpicture}[>=Stealth,x=2cm,y=2cm]
 \fill[manifoldblue!6] (0,0) rectangle (3,1.5);
 \fill[manifoldblue!13] (1,0) rectangle (2,1);
 \draw[manifoldblue!35] (0,0) grid[step=1] (3,1);
 \foreach \x in {0,1,2,3}
  {\draw[manifoldblue!35,dashed] (\x,1)--(\x,1.5);}
 \draw[->,manifoldred,very thick] (.3,.4)--(1.3,.4)
  node[midway,above=4pt] {$\widetilde\gamma$};
 \fill[manifoldblue!80!black] (.3,.4) circle(2pt) (1.3,.4) circle(2pt);
 \node[below=5pt] at (.3,.4) {$\widetilde p$};
 \node[below=5pt] at (1.3,.4) {$T_{(1,0)}\widetilde p$};
 \draw[->,thick,manifoldgold!85!black] (.3,1.175)--(1.3,1.175)
  node[midway,above=4pt] {$\rho(1,0)$};
 \node[align=left,anchor=west,font=\small,text width=4.2cm] at (3.35,.75)
  {$\widetilde T^2=\R^2$, $\mathrm{dev}=\id$.\\[5pt]
   $\rho(1,0)(x,y)=(x+1,y)$.\\[5pt]
   Końce czerwonej drogi mają ten sam obraz w $T^2$.};
\end{tikzpicture}
\caption{Rozwinięcie i holonomia płaskiego torusa.
Każdy kwadrat przedstawia jednostkową dziedzinę fundamentalną.
Podniesienie pętli nie zamyka się w modelu, lecz kończy
w punkcie przesuniętym przez holonomię.}
\label{fig:klein-rozwijanie-torusa}
\end{figure}

\begin{stwierdzenie}[Zupełność a globalny iloraz]
\label{prop:thur-dop-dev-kompletna}
Przy założeniach poprzedniego twierdzenia, jeśli $X$ jest zupełna jako rozmaitość
riemannowska, a atlas $(G,X)$ na spójnej $M$ wyznacza
zupełną metrykę lokalnie izometryczną z $X$, to
$\mathrm{dev}:\widetilde M\to X$ jest izometrią.
Wówczas $\rho$ jest injektywne, jego obraz jest dyskretny
i działa swobodnie, a $M$ jest ilorazem
$\rho(\pi_1(M))\backslash X$. W szczególności dotyczy to
zwartych rozmaitości z taką strukturą.
\end{stwierdzenie}
\begin{proof}
Metryka z $X$ skleja się na $M$, ponieważ zmiany
współrzędnych są izometriami. Odwzorowanie rozwijające
jest lokalną izometrią. Nakrycie uniwersalne z podniesioną
metryką jest zupełne: każda jego geodezyjna rzutuje się na
geodezyjną w $M$, którą przedłużamy na całą prostą i podnosimy
z powrotem przez nakrycie. Hopf--Rinow daje zupełność metryczną.

Pokażmy, że $F=\mathrm{dev}$ jest nakryciem całej $X$.
Dla $q_0=F(p_0)$ i dowolnego $q\in X$ wybieramy geodezyjną
od $q_0$ do $q$. Jej początkową prędkość podnosimy przez
$(\dd F_{p_0})^{-1}$ i rozwiązujemy równanie geodezyjne
w zupełnej $\widetilde M$. Lokalna izometryczność i jednoznaczność
geodezyjnych dają rzut równy wybranej geodezyjnej w $X$.
Jej koniec jest przeciwobrazem $q$, więc $F$ jest surjektywne.

Wybierzmy kulę normalną $B(q,r)$ w $X$. Dla każdego
$p\in F^{-1}(q)$ określamy na tej kuli
\[
 s_p(z)=\exp_p\bigl((\dd F_p)^{-1}\exp_q^{-1}(z)\bigr).
\]
Eksponenta w $p$ jest określona na całej przestrzeni stycznej
dzięki zupełności. Ponieważ lokalna izometria przenosi
geodezyjne z daną prędkością początkową, dla $|v|<r$ mamy
$F(\exp_p((\dd F_p)^{-1}v))=\exp_q(v)$.
Zatem $F\circ s_p=\id$.
Różniczkując, widzimy, że $\dd s_p$ jest izomorfizmem;
$s_p$ jest lokalnym dyfeomorfizmem i jest injektywne,
bo ma lewą odwrotność $F$. Jego obraz $U_p$ jest otwarty,
a $F|_{U_p}$ jest dyfeomorfizmem na $B(q,r)$.

Obrazy $U_p$ są parami rozłączne.
Gdyby $s_p(z)=s_{p'}(z)=w$, odwrócona radialna geodezyjna
od $z$ do $q$ miałaby dwa podniesienia od $w$.
Ich początkowe prędkości muszą być równe
$(\dd F_w)^{-1}$ zastosowanemu do tej samej prędkości,
więc jednoznaczność równania geodezyjnej daje $p=p'$.
Te obrazy wyczerpują $F^{-1}(B(q,r))$:
z dowolnego $w$ nad $z$ podnosimy tę radialną geodezyjną
za pomocą prędkości początkowej i zupełności.
Jej koniec $p$ leży nad $q$, a odwrócenie drogi daje $w=s_p(z)$.
Są to więc wszystkie arkusze nakrycia.
Nakrycie jednospójnej $X$ jest
trywialne, zatem $\mathrm{dev}$ jest globalną izometrią.
Równość~\eqref{eq:thur-dop-dev-equiv} identyfikuje
działanie grupy nakrycia z działaniem $\rho(\pi_1(M))$.
Ponieważ $\mathrm{dev}$ jest bijekcją, jądro $\rho$ jest
trywialne. Żaden nietrywialny element grupy nakrycia nie
ustala punktu, zatem dzięki równoważności to samo dotyczy
obrazu $\rho$.
Aby sprawdzić dyskretność w topologii $G$, wybierzmy małe
otoczenie $U$ punktu $\widetilde p$ rozłączne
z $T_\gamma U$ dla $\gamma\ne e$. Wtedy żaden
$\rho(\gamma)\ne e$ nie przenosi $F(\widetilde p)$ do
$F(U)$. Zbiór
$\{a\in G:a\cdot F(\widetilde p)\in F(U)\}$
jest otwartym otoczeniem jedności, którego przecięcie
z obrazem $\rho$ składa się tylko z jedności.
To jest właśnie dyskretność z
Definicji~\ref{def:klein-dyskretna-swobodna}.
Jeśli $M$ jest zwarta, jej każda metryka riemannowska
jest zupełna dzięki twierdzeniu Hopfa--Rinowa.
\end{proof}

\begin{przyklad}[Struktura niezupełna]
\label{prz:thur-dop-struktura-niezupelna}
Na otwartej kuli euklidesowej $B^n\subset\R^n$ wybierzmy
atlas otrzymany przez ograniczenie identyczności
z modelu $X=\R^n$ i grupy jego izometrii. Nakrycie
uniwersalne jest trywialne, $\mathrm{dev}:B^n\hookrightarrow\R^n$
jest inkluzją, a holonomia jest trywialna.
Metryka nie jest zupełna: odcinek zbliżający się do brzegu
ma skończoną długość.
Na przykład $p_j=(1-1/j,0,\ldots,0)$ dla $j\geq2$
jest ciągiem Cauchy'ego, ponieważ odcinek między jego punktami
leży w kuli i $d(p_j,p_k)=|1/j-1/k|$.
Nie ma granicy w $B^n$, gdyż jego jedyna granica euklidesowa
leży na brzegu. Dlatego nie dostajemy
ilorazu całego modelu $\R^n$.
\end{przyklad}

\begin{przyklad}[Rozwinięcie może nie być różnowartościowe]
\label{prz:klein-przebita-plaszczyzna}
Na $M=\R^2\setminus\{0\}$ bierzemy ograniczenie geometrii
euklidesowej. Nakrycie uniwersalne i rozwinięcie są dane przez
\[
 \widetilde M=(0,\infty)\times\R,\qquad
 \pi(r,t)=(r\cos t,r\sin t),\qquad
 \mathrm{dev}=\pi:\widetilde M\longrightarrow X=\R^2.
\]
Jest to nakrycie $M$, ponieważ w małym dysku omijającym zero
można wybrać jedną gałąź kąta, a inne różnią się o $2\pi k$.
Dziedzina jest jednospójna: $(r,t)\mapsto(\log r,t)$
jest dyfeomorfizmem na $\R^2$.
Przekształcenia nakrycia mają postać
$T_k(r,t)=(r,t+2\pi k)$.
Wektor pochodnej po $r$ jest $(\cos t,\sin t)$,
a po $t$ jest $(-r\sin t,r\cos t)$.
Ich iloczyny skalarne to $1,0,r^2$, więc metryka
podniesiona wynosi $\dd r^2+r^2\dd t^2$;
wyznacznik różniczki jest równy $r>0$.
Rozwinięcie jest zatem lokalną izometrią, lecz
$\mathrm{dev}(r,t+2\pi)=\mathrm{dev}(r,t)$
i jego obraz omija zero.
Holonomia jest trywialna, gdyż
$\mathrm{dev}\circ T_k=\mathrm{dev}$:
przy skutecznym działaniu izometrii jedynym elementem
ustalającym otwarty obraz jest jedność.
Mimo to $\pi_1(M)=\Z$ jest nietrywialna.
Nie ma sprzeczności z
Stwierdzeniem~\ref{prop:thur-dop-dev-kompletna}:
ciąg $(1/j,0)$ w $M$ jest Cauchy'ego i nie ma granicy w $M$.
Bez zupełności nie wolno więc wnioskować o injektywności
rozwinięcia ani reprezentacji holonomii.
\end{przyklad}

\par\smallskip
{\footnotesize\noindent\emph{Dalsza lektura:}
\href{https://users.monash.edu/~jpurcell/book/HypKnotTheory.pdf}
{J.\ S.\ Purcell, \emph{Hyperbolic Knot Theory}, \S\,3.1--3.3},
rozwinięcie, holonomia i zupełność struktur geometrycznych.\par}

\section{Przestrzeń euklidesowa i afiniczna}
\label{sec:klein-euklidesowa-afiniczna}

\begin{przyklad}[Geometria euklidesowa]
\label{prz:klein-euklidesowa}
\index{grupa euklidesowa}
Niech $n\geq1$. Grupa ruchów euklidesowych
\[
 E(n)=\R^n\rtimes\mathrm O(n)
 =\{(b,A):b\in\R^n,\ A\in\mathrm O(n)\}
\]
ma iloczyn $(b,A)(c,B)=(b+Ac,AB)$ i działa na $\R^n$
według $(b,A)\cdot x=b+Ax$. Współrzędne $(b,A)$ oraz
gładkość działań w czynnikach czynią z $E(n)$ grupę Liego.
Wzór iloczynu pochodzi ze składania ruchów:
$b+A(c+Bx)=(b+Ac)+(AB)x$.
Jednością jest $(0,I)$, a odwrotność wynosi
$(-A^{-1}b,A^{-1})$, bo
$b+A(-A^{-1}b)=0$ i $AA^{-1}=I$.
Jest to iloczyn półprosty z Rozdziału~5;
dla macierzy ortogonalnych $A^{-1}=A^{\mathsf T}$.
Translacje zapewniają przechodniość, a stabilizatorem $0$
jest $\{(0,A):A\in\mathrm O(n)\}\cong\mathrm O(n)$.
Twierdzenie~\ref{thm:klein-orbita-iloraz} daje
\[
 \R^n\cong E(n)/\mathrm O(n),\qquad
 \dim E(n)=n+\frac{n(n-1)}2,\qquad
 \dim\mathrm O(n)=\frac{n(n-1)}2.
\]
Działanie jest skuteczne: ruch ustalający $0$ ma $b=0$,
a ustalający następnie każdy wektor bazy ma $A=I$.
Macierz ortogonalna zachowuje $|x-y|$, więc zwykła metryka
euklidesowa jest $E(n)$-niezmiennicza. Są to wszystkie bijektywne
izometrie $\R^n$. Aby to udowodnić, odejmijmy obraz zera
od danej izometrii; tak otrzymana $f$ spełnia $f(0)=0$.
Z zachowania odległości mamy $|f(x)|=|x|$ oraz
\[
 \langle f(x),f(y)\rangle
 =\tfrac12\bigl(|f(x)|^2+|f(y)|^2-|f(x)-f(y)|^2\bigr)
 =\langle x,y\rangle.
\]
Stąd dla $w=f(x+y)-f(x)-f(y)$ zachodzi
$\langle w,f(z)\rangle
=\langle x+y,z\rangle-\langle x,z\rangle-\langle y,z\rangle=0$.
Surjektywność $f$ pozwala wybrać $z$ z $f(z)=w$,
więc $|w|^2=0$ i $f$ jest addytywna.
Addytywność daje $f(rx)=rf(x)$ dla $r\in\Q$,
a ciągłość przenosi równość na $r\in\R$.
Zatem $f$ jest liniowa. Jej macierz $A$ spełnia
$A^{\mathsf T}A=I$, bo $f$ zachowuje iloczyn skalarny.
Po przywróceniu odjętego obrazu zera dostajemy $x\mapsto b+Ax$.
Dopuszczamy odbicia;
wybór $\R^n\rtimes\mathrm{SO}(n)$ daje wariant zachowujący
orientację.
\end{przyklad}

\begin{przyklad}[Grupa afiniczna i przestrzeń afiniczna]
\label{prz:klein-afiniczna}
\index{grupa afiniczna}
\index{przestrzeń afiniczna}
Niech
\[
 \operatorname{Aff}(n,\R)
 =\R^n\rtimes\mathrm{GL}(n,\R),\qquad
 (b,A)\cdot x=b+Ax.
\]
Wzór na iloczyn jest taki sam jak dla $E(n)$; odwrotność
wynosi $(b,A)^{-1}=(-A^{-1}b,A^{-1})$. Ponieważ
$\mathrm{GL}(n,\R)$ jest otwarta w przestrzeni macierzy,
te wzory są gładkie. Stabilizatorem $0$ jest
$\mathrm{GL}(n,\R)$, działanie jest przechodnie dzięki
translacjom i skuteczne z tego samego powodu co wyżej.
Zatem
\[
 \R^n\cong\operatorname{Aff}(n,\R)/\mathrm{GL}(n,\R),
 \qquad (n+n^2)-n^2=n.
\]
Jako rozmaitość jest to znowu $\R^n$, lecz większa grupa
zachowuje inne własności. Odwzorowanie $x\mapsto Ax+b$
wysyła proste na proste i zachowuje podział odcinka w danym
stosunku, gdyż
$(1-t)f(x)+tf(y)=f((1-t)x+ty)$. Nie musi zachować długości:
$x\mapsto2x$ należy do grupy afinicznej.
Właściwie żadna dodatnio określona metryka riemannowska
na $\R^n$ nie jest niezmiennicza względem całej tej grupy.
Jej stabilizator zera zawiera $x\mapsto2x$; niezmienniczość
w zerze wymagałaby $g_0(2v,2v)=g_0(v,v)$ dla każdego $v$,
co przeczy dodatniej określoności:
dwuliniowość dawałaby $4g_0(v,v)=g_0(v,v)$,
czyli $g_0(v,v)=0$ również dla $v\ne0$.
\end{przyklad}

\section{Przestrzeń rzutowa i sfera}
\label{sec:klein-rzutowa-sfera}

Przestrzeń $\mathbb{RP}^n$ z Przykładu~\ref{prz:RPn}
składa się z jednowymiarowych podprzestrzeni $\R^{n+1}$;
punkt zapisujemy $[x_0:\cdots:x_n]$. W tej geometrii istotne
są proste rzutowe i ich wzajemne położenie, a nie odległości.

\begin{przyklad}[Rzutowa grupa liniowa]
\label{prz:klein-rzutowa}
\index{rzutowa grupa liniowa}
Niech $n\geq1$ i
\[
 \mathrm{PGL}(n+1,\R)
 =\mathrm{GL}(n+1,\R)/(\R_*I),
 \qquad \R_*=\R\setminus\{0\}.
\]
Podgrupa skalarnych macierzy $\R_*I$ jest domknięta
w $\mathrm{GL}(n+1,\R)$, więc iloraz jest grupą Liego.
Klasa macierzy $A$ działa na $\mathbb{RP}^n$ przez
$[A]\cdot[x]=[Ax]$; zmiana $A$ przez niezerowy skalar
nie zmienia obrazu. W mapach rzutowych z
Przykładu~\ref{prz:RPn} działanie ma współrzędne będące
ilorazami funkcji liniowych o niezerowym mianowniku,
więc jest gładkie. Grupa $\mathrm{GL}(n+1,\R)$ przenosi
dowolną prostą liniową na dowolną inną, zatem działanie
jest przechodnie.

Wybierzmy $p=[1:0:\cdots:0]$. Jego stabilizator $P$
składa się z klas blokowych macierzy
\[
 \begin{pmatrix}\lambda&\beta\\0&B\end{pmatrix},
 \qquad \lambda\in\R_*,\quad
 \beta\in\R^{1\times n},\quad B\in\mathrm{GL}(n,\R).
\]
Warunek zerowego lewego dolnego bloku mówi dokładnie,
że pierwsza kolumna jest wielokrotnością $e_0$.
Stąd $\mathbb{RP}^n\cong\mathrm{PGL}(n+1,\R)/P$.
Kontrola wymiarów daje
\[
 \dim\mathrm{PGL}(n+1,\R)=(n+1)^2-1,
 \quad \dim P=n^2+n,
 \quad \dim\mathbb{RP}^n=n.
\]
Wymiar $P$ policzyliśmy jako
$1+n+n^2-1$: ostatnie $-1$ usuwa skalarne macierze.

Działanie $\mathrm{PGL}$ jest skuteczne. Jeśli $A$ zachowuje
każdą prostą, to $Ae_i=\lambda_i e_i$ dla wektorów bazy.
Zachowanie prostej $[e_i+e_j]$ wymusza
$\lambda_i=\lambda_j$, a więc $A$ jest skalarna i
$[A]$ jest jednością w $\mathrm{PGL}$.

Pełna geometria rzutowa nie ma niezmienniczej metryki
riemannowskiej. W mapie $[1:x_1:\cdots:x_n]$ element
$[\operatorname{diag}(1,2I_n)]\in P$ działa jako
$x\mapsto2x$. Różniczka w $p$ mnoży każdy wektor przez
$2$, co przeczy niezmienniczości dodatnio określonego
iloczynu skalarnego w $T_p\mathbb{RP}^n$. Można
\emph{wybrać} na $\mathbb{RP}^n$ metrykę pochodzącą ze
sfery, lecz jej grupa symetrii jest mniejsza niż pełne
$\mathrm{PGL}(n+1,\R)$.
\end{przyklad}

\begin{przyklad}[Sfera jako przestrzeń obrotów]
\label{prz:klein-sfera}
\index{sfera!przestrzeń jednorodna}
Dla $n\geq1$ grupa $\mathrm{SO}(n+1)$ działa na
$S^n\subset\R^{n+1}$ przez obroty. Działanie jest
przechodnie: jednostkowe wektory $x,y$ można uzupełnić
do dodatnio zorientowanych baz ortonormalnych, a macierz
przeprowadzająca jedną bazę w drugą leży w
$\mathrm{SO}(n+1)$. Stabilizator $e_0$ składa się z
$\operatorname{diag}(1,B)$, $B\in\mathrm{SO}(n)$.
Istotnie, jeśli $Ae_0=e_0$, to dla $v\perp e_0$
mamy $\langle Av,e_0\rangle
=\langle Av,Ae_0\rangle=\langle v,e_0\rangle=0$.
Macierz zachowuje więc rozkład
$\R e_0\oplus e_0^\perp$; jej blok na $e_0^\perp$
jest ortogonalny i ma wyznacznik $\det A=1$.
Zatem
\[
 S^n\cong\mathrm{SO}(n+1)/\mathrm{SO}(n),
 \qquad
 \frac{n(n+1)}2-\frac{n(n-1)}2=n.
\]
Dla $n=1$ konwencja $\mathrm{SO}(1)=\{1\}$ sprawia,
że wzór mówi $S^1\cong\mathrm{SO}(2)$; nie ma tu
jednowymiarowego stabilizatora. Działanie jest skuteczne,
bo macierz ustalająca wszystkie punkty sfery ustala
każdy wektor przestrzeni. Metryka kulista otrzymana
przez ograniczenie metryki euklidesowej jest zachowana
przez $\mathrm{SO}(n+1)$.

Antypody $x$ i $-x$ na sferze wyznaczają tę samą
prostą liniową, stąd $S^n/\{x\sim-x\}\cong\mathbb{RP}^n$.
Metryka kulista schodzi na ten iloraz, ponieważ
antypodalność jest izometrią. Ta metryka jest dodatkową
strukturą: nie jest niezmiennicza względem pełnej grupy
rzutowej z poprzedniego przykładu. Grupa
$\mathrm O(n+1)$ działa na $\mathbb{RP}^n$ przez izometrie
tej metryki, z jądrem $\{I,-I\}$; jej skuteczny obraz
to $\mathrm{PO}(n+1)=\mathrm O(n+1)/\{I,-I\}$.
Jądro ustaliliśmy już w poprzednim przykładzie:
macierz zachowująca każdą prostą jest skalarna,
a ortogonalność skalarnej $\lambda I$ wymaga $\lambda=\pm1$.
\end{przyklad}

\section{Torus: iloraz, grupa i geometria płaska}
\label{sec:klein-torus}

\begin{przyklad}[Dwa działania przechodnie na torusie]
\label{prz:klein-torus}
\index{torus!geometria Kleina}
Niech $n\geq1$ i $T^n=\R^n/\Z^n$, jak w
Przykładzie~\ref{prz:torus}. Rozdział~\ref{ch:grupy-liego}
wyposażył tę rozmaitość w strukturę grupy Liego:
$[x]+[y]=[x+y]$. Dlatego pierwszym, skutecznym modelem
jest działanie grupy $G=T^n$ na $X=T^n$ przez translacje.
Stabilizator $[0]$ to $H=\{[0]\}$ i
\[
 T^n\cong T^n/\{[0]\}.
\]

Możemy zamiast tego pozostawić grupę $G=(\R^n,+)$ i
zdefiniować $a\cdot[x]=[a+x]$. Działanie jest gładkie
i przechodnie, a stabilizator $[0]$ jest równy
$H=\Z^n$. Otrzymujemy zatem ten sam torus jako
\[
 T^n\cong\R^n/\Z^n.
\]
Tym razem działanie \emph{nie jest skuteczne}: każdy
$a\in\Z^n$ działa na każdym $[x]$ jak jedność.
Jądro~\eqref{eq:klein-jadro} wynosi właśnie $\Z^n$,
a skuteczna grupa ilorazowa to $\R^n/\Z^n=T^n$.
Podgrupa $\Z^n$ jest dyskretna, więc wzór na wymiar
ma postać $n-0=n$.

Zwykła metryka na $\R^n$ jest niezmiennicza względem
translacji przez $\Z^n$, toteż przechodzi na płaską
metrykę torusa. Dokładniej, dla dowolnej symetrycznej
dodatnio określonej macierzy $B$ metryka
$g_B=\sum_{i,j}B_{ij}\dd x^i\dd x^j$ jest stała,
a różniczka translacji jest identycznością,
więc wszystkie translacje zachowują $g_B$.
Na przykład weźmy na tym samym $T^2$ pętle
$\alpha(s)=[(s,0)]$ i $\beta(s)=[(0,s)]$, $s\in[0,1]$.
Ich prędkości w mapach ilorazowych to odpowiednio
$e_1$ i $e_2$, zatem normy wynoszą $\sqrt{B_{11}}$
i $\sqrt{B_{22}}$.
Dla $B=I_2$ i $B=\operatorname{diag}(1,4)$ długość
$\alpha$ jest więc równa $\int_0^1 1\,\dd s=1$,
natomiast długość $\beta$ wynosi odpowiednio $1$
i $\int_0^1 2\,\dd s=2$.
Metryka zależy więc od wybranego iloczynu
skalarnego na $\R^n$; samo przedstawienie torusa jako
grupy Liego nie wyznacza długości. Torus obrotowy
zanurzony w $\R^3$ z Przykładu~\ref{prz:torus-liego}
ma z kolei metrykę indukowaną przez to zanurzenie.
Dla parametryzacji
$F(\theta,\phi)=((R+r\cos\theta)\cos\phi,
(R+r\cos\theta)\sin\phi,r\sin\theta)$, gdzie $R>r>0$, mamy
\begin{align*}
 F_\theta&=(-r\sin\theta\cos\phi,-r\sin\theta\sin\phi,r\cos\theta),\\
 F_\phi&=(-(R+r\cos\theta)\sin\phi,(R+r\cos\theta)\cos\phi,0).
\end{align*}
Iloczyn mieszany znika, bo jego dwa wyrazy zawierające
$\sin\phi\cos\phi$ mają przeciwne znaki.
Z $\sin^2\phi+\cos^2\phi=1$ otrzymujemy
$|F_\theta|^2=r^2(\sin^2\theta+\cos^2\theta)=r^2$
i $|F_\phi|^2=(R+r\cos\theta)^2$.
Wzór~\eqref{eq:metryka-pullback} daje więc
$F^*g_{\R^3}=r^2\dd\theta^2+(R+r\cos\theta)^2\dd\phi^2$.
Jej zależność od $\theta$ pokazuje, że translacje współrzędnej
$\theta$ nie są na ogół izometriami.
Trzeba więc odróżniać
rozmaitość, jej strukturę grupową i wybraną metrykę.
\end{przyklad}

\section{Dysk Poincarégo jako iloraz grupy Liego}
\label{sec:klein-dysk}

\begin{przyklad}[Działanie $\mathrm{PSL}(2,\R)$]
\label{prz:klein-dysk}
\index{dysk Poincarégo}
Zacznijmy od górnej półpłaszczyzny
$\mathbb H^2=\{z=x+iy\in\C:y>0\}$. Niech
\[
 G=\mathrm{PSL}(2,\R)
 =\mathrm{SL}(2,\R)/\{I,-I\}.
\]
Podgrupa $\{I,-I\}$ jest domknięta i normalna,
więc $G$ jest grupą Liego. Jej wymiar wynosi trzy:
$\mathrm{SL}(2,\R)=\det^{-1}(1)$ jest regularną
hiperpowierzchnią w czterowymiarowej przestrzeni
macierzy, bo $\dd(\det)_A(A)=2\det A=2$;
obliczamy tę pochodną z
$\det(A+sA)=(1+s)^2\det A$, różniczkując w $s=0$.
W szczególności różniczka wyznacznika nie jest zerowa.
Zatem
iloraz przez dwuelementową podgrupę
nie zmienia wymiaru. Grupa $G$ działa na
$\mathbb H^2$ przez przekształcenia ułamkowo liniowe
\[
 \left[\begin{pmatrix}a&b\\c&d\end{pmatrix}\right]
 \cdot z=\frac{az+b}{cz+d},\qquad ad-bc=1.
\]
Znak całej macierzy nie zmienia ułamka. Mianownik
nie znika: jeśli $c\ne0$, jego zero $-d/c$ byłoby rzeczywiste,
a jeśli $c=0$, warunek wyznacznika wymusza $d\ne0$.
Oznaczmy ułamek przez $f_A(z)$.
Dla $B=\begin{psmallmatrix}\alpha&\beta\\\gamma&\delta\end{psmallmatrix}$
podstawienie jednego ułamka w drugi daje
\[
 f_A(f_B(z))
 =\frac{(a\alpha+b\gamma)z+(a\beta+b\delta)}
        {(c\alpha+d\gamma)z+(c\beta+d\delta)}
 =f_{AB}(z).
\]
Ponadto $f_I(z)=z$, a więc po sprawdzeniu, że obraz leży
w półpłaszczyźnie, otrzymamy działanie grupy z odwrotnością
$f_{A^{-1}}$.
W tym celu mnożymy licznik i mianownik przez
$c\bar z+d$. Licznik staje się
$ac|z|^2+bd+ad\,z+bc\,\bar z$ i ma część urojoną
$(ad-bc)y=y$. Reguła różniczkowania ilorazu daje natomiast
$f_A'(z)=(a(cz+d)-c(az+b))/(cz+d)^2$.
Stąd
\begin{equation}\label{eq:klein-mobius-pochodna}
 \Im\frac{az+b}{cz+d}=\frac{y}{|cz+d|^2},
 \qquad
 \frac{\dd}{\dd z}\frac{az+b}{cz+d}
 =\frac1{(cz+d)^2}.
\end{equation}
Pierwsza równość pokazuje, że obraz pozostaje w
$\mathbb H^2$; druga, że działanie jest gładkie.

Działanie jest przechodnie. Dla $z=x+iy$ macierz
\[
 A_{x,y}=\begin{pmatrix}
 \sqrt y&x/\sqrt y\\0&1/\sqrt y
 \end{pmatrix}\in\mathrm{SL}(2,\R)
\]
ma wyznacznik $\sqrt y/\sqrt y=1$ i wysyła
$i$ w $(\sqrt y\,i+x/\sqrt y)/(1/\sqrt y)=x+iy$.
Stabilizator $i$ obliczamy z
$(ai+b)/(ci+d)=i$: po porównaniu części rzeczywistej
i urojonej otrzymujemy $b=-c$ oraz $a=d$.
Warunek wyznacznika jest wtedy $a^2+c^2=1$.
Zatem stabilizator jest \emph{obrazem} $\mathrm{SO}(2)$
w $\mathrm{PSL}(2,\R)$, czyli
\[
 K=\mathrm{SO}(2)/\{I,-I\},\qquad
 \mathbb H^2\cong\mathrm{PSL}(2,\R)/K,
 \qquad 3-1=2.
\]
Grupa $K$ jest izomorficzna z kołem obrotów,
lecz nie jest dosłownie podgrupą macierzy
$\mathrm{SO}(2)\subset\mathrm{SL}(2,\R)$: w $G$
utożsamiliśmy macierze $I$ i $-I$. Konkretny izomorfizm
na zwykłą grupę obrotów ma postać $[R_\theta]\mapsto R_{2\theta}$.
Jest dobrze określony, gdyż $R_{\theta+\pi}=-R_\theta$,
a $R_{2(\theta+\pi)}=R_{2\theta}$.
Dodawanie kątów dowodzi zachowania iloczynu;
jądro składa się z $\theta\in\pi\Z$, czyli jednej klasy w $K$,
a każdy obrót $R_\alpha$ jest obrazem $[R_{\alpha/2}]$.

Metryka hiperboliczna
$g_{\mathbb H}=((\dd x)^2+(\dd y)^2)/y^2$ z
Równania~\eqref{eq:hip-polprzestrzen} jest zachowana
przez $G$. Istotnie, przekształcenie holomorficzne
w~\eqref{eq:klein-mobius-pochodna} mnoży euklidesowy
element długości przez $|cz+d|^{-2}$, zaś wysokość
$y$ przez ten sam czynnik. W ilorazie definiującym
$g_{\mathbb H}$ czynniki te się skracają. W zapisie tensorowym
jest to rachunek
\[
 f_A^*g_{\mathbb H}
 =\frac{|f_A'(z)|^2|\dd z|^2}{(\Im f_A(z))^2}
 =\frac{|cz+d|^{-4}|\dd z|^2}{y^2|cz+d|^{-4}}
 =g_{\mathbb H}.
\]
Jeśli element $G$ ustala wszystkie $z\in\mathbb H^2$,
to wielomian
$cz^2+(d-a)z-b$ znika na zbiorze otwartym, więc
$c=b=0$, $d=a$ i $a=\pm1$. W $G$ jest to jedność:
działanie jest skuteczne.

Przekształcenie Cayleya
\[
 C:\mathbb H^2\longrightarrow D,
 \qquad C(z)=\frac{z-i}{z+i},\qquad
 D=\{w\in\C:|w|<1\},
\]
ma $C(i)=0$. Rozwiązując $w(z+i)=z-i$ względem $z$,
dostajemy $C^{-1}(w)=i(1+w)/(1-w)$.
Dla $|w|<1$ mianownik nie znika, a
$\Im C^{-1}(w)=(1-|w|^2)/|1-w|^2>0$:
otrzymujemy tę równość, mnożąc licznik i mianownik
przez $1-\bar w$ i biorąc część urojoną.
Z kolei dla $z=x+iy$ różnica
$|z+i|^2-|z-i|^2=(y+1)^2-(y-1)^2=4y$ jest dodatnia.
Zatem $C$ i $C^{-1}$ mają wskazane dziedziny i są gładkie.
To dowodzi, że $C$ jest dyfeomorfizmem.
Ten sam rachunek różnicy kwadratów oraz
$C'(z)=((z+i)-(z-i))/(z+i)^2=2i/(z+i)^2$ dają
\[
 1-|C(z)|^2=\frac{4y}{|z+i|^2},
 \qquad |C'(z)|=\frac2{|z+i|^2}
\]
wynika przez podstawienie
\[
 C^*\!\left(\frac{4|\dd w|^2}{(1-|w|^2)^2}\right)
 =\frac{4\cdot4|z+i|^{-4}}{16y^2|z+i|^{-4}}|\dd z|^2
 =\frac{|\dd z|^2}{y^2}=g_{\mathbb H}.
\]
Tak otrzymujemy model dysku Poincarégo z
Równania~\eqref{eq:hip-kula} dla $n=2$ oraz jego
opis $D\cong\mathrm{PSL}(2,\R)/K$.
Po przeniesieniu działania przez $C$ stabilizator
zera działa obrotami dysku.
Przyjmujemy $R_\theta=
 \begin{psmallmatrix}\cos\theta&-\sin\theta\\
\sin\theta&\cos\theta\end{psmallmatrix}$.
Po oznaczeniu $u=\cos\theta$, $v=\sin\theta$ mamy
\[
 C(R_\theta\cdot z)
 =\frac{(u-iv)z-(v+iu)}{(u+iv)z-v+iu}
 =\frac{(u-iv)(z-i)}{(u+iv)(z+i)}
 =e^{-2i\theta}C(z).
\]
Minus w kącie wynika z wybranego wzoru na $C$
i standardowej macierzy $R_\theta$.
Nie zmienia wcześniejszego abstrakcyjnego izomorfizmu
$[R_\theta]\mapsto R_{2\theta}$.

Brzeg $|w|=1$ nie należy do przestrzeni.
Dla promieniowej drogi od $0$ do punktu o module $r<1$
długość wynosi
\[
 L(r)=\int_0^r\frac{2\,\dd s}{1-s^2}
 =\int_0^r\left(\frac1{1-s}+\frac1{1+s}\right)\dd s
 =\log\frac{1+r}{1-r}\longrightarrow\infty
 \quad(r\to1^-).
\]
Wartość całki otrzymujemy z funkcji pierwotnej
$-\log(1-s)+\log(1+s)$ i wartości w $s=0$.
Każda inna kawałkami gładka droga $w(t)$, $0\leq t<T$,
której moduł $r(t)=|w(t)|$ dąży do $1$, także ma
nieskończoną długość. Istotnie, $|r'(t)|\leq|w'(t)|$
prawie wszędzie, zatem dla $t_0<t<T$ mamy
\[
 \int_{t_0}^{t}\frac{2|w'(s)|}{1-|w(s)|^2}\,\dd s
 \geq\left|\int_{t_0}^{t}\frac{2r'(s)}{1-r(s)^2}\,\dd s\right|
 =|L(r(t))-L(r(t_0))|\longrightarrow\infty.
\]
\end{przyklad}

\begin{figure}[H]
\centering
\begin{tikzpicture}[>=Stealth,scale=1.05]
 \fill[manifoldblue!6] (-1.5,0) rectangle (1.5,3);
 \draw[manifoldblue!45] (-1.6,0)--(1.6,0);
 \draw[->,very thick,manifoldred] (0,.08)--(0,2.9);
 \foreach \yy/\txt in {.5/{i/2},1/{i},2/{2i}}
  {\fill[manifoldblue!80!black] (0,\yy) circle(2pt);
   \node[left=6pt,font=\small] at (0,\yy) {$\txt$};}
 \node at (0,3.35) {$\mathbb H^2$};
 \draw[->,thick,manifoldgold!85!black] (2,1.4)--(4.2,1.4)
  node[midway,above=5pt] {$C(z)=\dfrac{z-i}{z+i}$};
 \begin{scope}[shift={(6.2,1.5)}]
  \fill[manifoldblue!8] (0,0) circle(1.3);
  \draw[manifoldblue!70] (0,0) circle(1.3);
  \draw[->,very thick,manifoldred] (-1.23,0)--(1.23,0);
  \foreach \xx/\txt in {-.433333/{-1/3},.433333/{1/3}}
   {\fill[manifoldblue!80!black] (\xx,0) circle(2pt);
    \node[above=6pt,font=\small] at (\xx,0) {$\txt$};}
  \fill[manifoldblue!80!black] (0,0) circle(2pt);
  \node[below=6pt,font=\small] at (0,0) {$0$};
  \draw[draw=manifoldblue!70,fill=white,thick] (-1.3,0) circle(2.5pt) (1.3,0) circle(2.5pt);
  \node[left=4pt,font=\small] at (-1.3,0) {$-1$};
  \node[right=4pt,font=\small] at (1.3,0) {$1$};
 \end{scope}
 \node at (6.2,3.35) {$D$};
\end{tikzpicture}
\caption{Izometria Cayleya przenosi oś $z=iy$, $y>0$,
na średnicę $(-1,1)$: $C(iy)=(y-1)/(y+1)$.
Punkty $i/2,i,2i$ przechodzą odpowiednio w $-1/3,0,1/3$.
Strzałki wskazują rosnące $y$; końce średnicy leżą poza dyskiem.
Położenia na kartce są euklidesowe, a odległości mierzymy
metrykami hiperbolicznymi.}
\label{fig:klein-cayley-os-srednica}
\end{figure}

\section{Jak przechodzi się między geometriami}
\label{sec:klein-przejscia}

Grupy i ilorazy porządkują opisane przykłady, lecz
przejścia między nimi mają różny sens.

Po pierwsze, na \emph{tej samej} przestrzeni $\R^n$ mamy
inkluzję grup przekształceń
\[
 E(n)=\R^n\rtimes\mathrm O(n)
 \ \subset\ \operatorname{Aff}(n,\R)
 =\R^n\rtimes\mathrm{GL}(n,\R).
\]
Większa grupa oznacza mniej własności zachowywanych przez
wszystkie ruchy: długość jest euklidesowa, ale nie afiniczna;
współliniowość jest afiniczna. Oba modele mają tę samą
rozmaitość punktów, choć inne stabilizatory.

Po drugie, odwzorowanie $x\mapsto[1:x_1:\cdots:x_n]$
identyfikuje $\R^n$ z otwartym obszarem afinicznym
$U_0=\{[x_0:\cdots:x_n]:x_0\ne0\}$ w $\mathbb{RP}^n$.
Homomorfizm
\[
 (b,A)\longmapsto
 \left[\begin{pmatrix}1&0\\b&A\end{pmatrix}\right]
 \quad\text{z}\quad
 \operatorname{Aff}(n,\R)
 \longrightarrow\mathrm{PGL}(n+1,\R)
\]
jest różnowartościowy: równość dwóch klas wymagałaby
skalarnej proporcjonalności macierzy, a ich lewe górne
wyrazy równe $1$ wymuszają skalar $1$.
Zachowanie iloczynu sprawdzamy mnożąc bloki:
\[
 \begin{pmatrix}1&0\\b&A\end{pmatrix}
 \begin{pmatrix}1&0\\c&B\end{pmatrix}
 =\begin{pmatrix}1&0\\b+Ac&AB\end{pmatrix}.
\]
Realizuje zwykłe działanie
$x\mapsto Ax+b$ na $U_0$, ponieważ
$\begin{psmallmatrix}1&0\\b&A\end{psmallmatrix}
\binom{1}{x}=\binom{1}{b+Ax}$.
Jej obraz jest dokładnie stabilizatorem hiperpłaszczyzny
$x_0=0$: macierz zachowująca tę hiperpłaszczyznę ma
pierwszy wiersz $(\lambda,0,\ldots,0)$ z $\lambda\ne0$,
a podzielenie całej macierzy przez $\lambda$ daje powyższy blok.
Pełna grupa rzutowa
zawiera też przekształcenia, które nie zachowują
wybranego obszaru afinicznego: punkt z $\R^n$ może przejść
na hiperpłaszczyznę w nieskończoności. Zmienia się więc
zarówno grupa, jak i globalna przestrzeń.
Na przykład macierz zamieniająca współrzędne $x_0$ i $x_1$
wysyła punkt $[1:0:\cdots:0]$ z $U_0$ w
$[0:1:0:\cdots:0]$, który leży na tej hiperpłaszczyźnie.

Po trzecie, ilorazy zmieniają \emph{topologię} modelu.
Torus płaski powstaje z $\R^n$ przez identyfikację
$x\sim x+m$, $m\in\Z^n$, a przestrzeń rzutowa ze
sfery przez identyfikację punktów antypodalnych.
W obu przypadkach należy odróżnić grupę działającą na
przestrzeni pokrywającej od skutecznej grupy działającej
na ilorazie. Model dysku i półpłaszczyzny jest jeszcze
innym przejściem: przekształcenie Cayleya jest
izometrią między dwiema postaciami \emph{tej samej}
geometrii hiperbolicznej.

Geometrie trójwymiarowe omawiane przy twierdzeniu
Thurstona również można zapisywać jako przestrzenie
jednorodne grup Liego, lecz sama przechodniość działania
nie wystarcza do ich klasyfikacji: w tamtym zagadnieniu
ważne są także wybrana metryka i istnienie odpowiednich
ilorazów przez dyskretne grupy.

\par\smallskip
{\footnotesize\noindent\emph{Dalsza lektura:}
\href{https://math.ucr.edu/home/baez/erlangen/Klein-erlangen.pdf}
{F.\ Klein, \emph{A Comparative Review of Recent Researches in Geometry}}
(oryginalny program erlangeński z 1872 roku, przekład);
\href{https://link.springer.com/book/10.1007/978-1-4419-9982-5}
{J.\ M.\ Lee, \emph{Introduction to Smooth Manifolds}, wyd.~2},
rozdziały~20--21 (podgrupy domknięte, działania i ilorazy);
\href{https://u.math.biu.ac.il/~solomyb/TEACH/21/GrAct/Lec5b.pdf}
{B.\ Solomyak, \emph{The Hyperbolic Plane}},
wykład o działaniu $\mathrm{PSL}(2,\R)$ i stabilizatorze $i$.\par}
